\documentclass[11pt]{article}

\usepackage[T1]{fontenc}
\usepackage[utf8]{inputenc}
\usepackage{lmodern}
\usepackage[margin=1.05in]{geometry}
\usepackage{amsmath,amssymb,amsthm,mathtools}
\usepackage{booktabs}
\usepackage{array}
\usepackage{enumitem}
\usepackage{xcolor}
\usepackage{hyperref}
\usepackage{xurl}
\DeclareUnicodeCharacter{2082}{\ensuremath{_2}}
\DeclareUnicodeCharacter{2083}{\ensuremath{_3}}
\DeclareUnicodeCharacter{2084}{\ensuremath{_4}}

\hypersetup{
  colorlinks=true,
  linkcolor=blue!55!black,
  citecolor=green!35!black,
  urlcolor=blue!55!black
}
\setlist{itemsep=2pt,topsep=4pt}
\setlength{\parindent}{0pt}
\setlength{\parskip}{5pt}
\setlength{\emergencystretch}{3em}

\newtheorem{theorem}{Theorem}[section]
\newtheorem{proposition}[theorem]{Proposition}
\newtheorem{lemma}[theorem]{Lemma}
\newtheorem{definition}[theorem]{Definition}
\newtheorem{remark}[theorem]{Remark}

\newcommand{\R}{\mathbb R}
\newcommand{\C}{\mathbb C}
\newcommand{\Q}{\mathbb Q}
\newcommand{\Z}{\mathbb Z}
\newcommand{\N}{\mathbb N}
\newcommand{\HGAS}{\mathrm{HGAS}}
\newcommand{\dd}{\,\mathrm d}
\newcommand{\Lean}{\textsc{Lean}}
\newcommand{\Mathlib}{\textsc{Mathlib}}
\newcommand{\statusproved}{\textcolor{green!40!black}{\textbf{proved}}}
\newcommand{\statusinterface}{\textcolor{orange!70!black}{\textbf{interface}}}
\newcommand{\statusopen}{\textcolor{red!65!black}{\textbf{open}}}
\newcommand{\leanname}[1]{\texttt{\detokenize{#1}}}
\newcommand{\leanmodule}[1]{\path{#1}}

\title{\textbf{Homogeneous Polynomial GAS and Computability}\\
  A Human-Readable Companion to the Lean Development}
\author{Formalization project companion}
\date{24 July 2026}

\begin{document}
\maketitle

\begin{center}
\fbox{\begin{minipage}{0.91\textwidth}
\textbf{Historical snapshot; superseded.}
This document records an intermediate formalization state on 24 July 2026
and intentionally preserves obligations that were open at that time.  It
must not be used as the current theorem-status report.  The completed,
audited proof companion is
\path{lean/docs/hgas_sigma1_complete_formalization_companion.tex}.
\end{minipage}}
\end{center}

\begin{abstract}
This document explains the mathematical proof architecture implemented in
the \Lean{} project \leanmodule{lean/HGAS}.  It is written for a reader with
a graduate background in analysis, dynamical systems, algebra, and basic
computability theory.  Its purpose is twofold: to make the formal proof
readable without requiring familiarity with \Lean{}, and to identify the
precise boundary between kernel-checked results and the integration theorems
that remain to be formalized.

The development already proves, without project-specific axioms, the exact
definition and angular characterization of homogeneous global asymptotic
stability, the computably enumerable upper bound in every fixed odd degree,
an unconditional primitive-recursive MRDP-to-fixed-DQL source, the required
clock and invariant-measure infrastructure, an ergodic-maximizer theorem for
continuous real flows on compact metric spaces, a substantial modular and
Jacobi theta theory (including the second-order Delta identity, global
projective saturation, an effective ten-coordinate polynomial compiler,
arbitrary finite-family complex realification, and an effective linked
unary--Jacobi compiler on twenty-eight real coordinates), exact conversion
between physical angular velocity in radians and additive-circle frequency
in turns, an analytic coupled controller--phase flow with exact ODE
semantics, and a complete ergodic classification of its clock marginals:
the controller tuple is almost surely locked, the phase projection is
semiconjugate to a constant rotation on that full-measure fibre, and the
phase marginal is Haar after conversion to turns.
On the finite-syntax side it now also proves a dynamically coupled
two-controller linked-Jacobi compiler on thirty-six real coordinates, with
controller/target bidegree at most \((3,8)\), ordinary total degree at most
\(9\), and two literal fifty-eight-dimensional compact-pipeline instances.
The provisional all-affine layout uses clearing exponent \(10\); the
geometry-faithful mixed layout retains the four controller and phase
circles as compact blocks, centrally compactifies only the linked block,
and therefore uses the sharp affine clearing exponent \(8\), with global
tangency even at the affine equator.  Further proved pieces include a
primitive-recursive right-hand-side oscillator compiler, an exact
classification of ergodic invariant probabilities for the corresponding
resonant additive-torus oscillator as Haar measures on unique invariant
graphs, a physical-clock detector-average theorem, conditional
recurrent-measure classification along any supplied modular factor map, the
algebraic and ODE portions of the compact accumulator compiler,
component-restricted finite invariant graphs, the homogeneous product-sphere
wrapper, effective odd-degree-to-cubic reduction, and effective cofinal
dimension/degree propagation.

The full lower-bound theorem is \emph{not yet claimed by the Lean
development}.  Three central integration tasks remain visible: the
Dani--Smillie/Hubbard--Miller horocycle classification in the exact quotient
used here; the recurrent modular/Jacobi selector and its state-map and
identification with the proved finite-skew-product interface (or a new
nontrivial-cover lifting theorem); and composition of the now separate
controller, phase, oscillator, detector, and linked-Jacobi layers into one
canonical fixed-resource source family.  The last task still includes the
source-specific invariant component and primitive for its positive graph,
the final-layout tangency proof, and the ergodic nonzero-drift alternative.
The coupled-controller Haar theorem itself is now complete; what remains is
to identify the appropriate factor of the complete enlarged polynomial
source with that analytic clock flow and with the detector/oscillator
models.  The formal files represent the remaining inputs as propositions
or certificate fields, never as axioms.
\end{abstract}

\tableofcontents

\section{How to read the status of this document}

There are three kinds of statements below.

\begin{center}
\begin{tabular}{>{\raggedright\arraybackslash}p{0.17\textwidth}
                >{\raggedright\arraybackslash}p{0.74\textwidth}}
\toprule
\statusproved &
A theorem whose proof is checked by the Lean kernel and whose source contains
no \leanname{sorry}, \leanname{admit}, project \leanname{axiom}, or
\leanname{unsafe} declaration.\\
\statusinterface &
A proved implication or a structure definition whose hypotheses deliberately
name the remaining mathematical input.  Such a declaration is useful, but it
does not prove that its hypotheses hold.\\
\statusopen &
A mathematical theorem or piece of source assembly still absent from the
kernel-checked dependency chain.\\
\bottomrule
\end{tabular}
\end{center}

This distinction is essential.  For example,
\leanname{CanonicalCompactAccumulatorGap} is a well-designed structure with
two fields, but the existence of an inhabitant for the universal arithmetic
source is still open.  Similarly,
\leanname{GammaTwoDaniSmillieTheorem} is the \emph{definition of a
proposition}; it is not an asserted axiom and not yet a proved theorem.
The theorems which take an argument
\leanname{hDS : GammaTwoDaniSmillieTheorem} are therefore proved conditional
implications, not proofs of that argument.

At the time of this companion, every module described as proved has been
built with warnings treated as errors; the integrated public surface is
\leanmodule{lean/HGAS.lean}, while newly stabilized modules are also named
explicitly below.  The intended final trust criterion is stronger than
compilation:
the headline theorem should depend only on the ordinary logical foundations
reported by \Lean{}/\Mathlib{} (typically \leanname{propext},
\leanname{Classical.choice}, and \leanname{Quot.sound}), and on no
project-defined postulate.

\subsection{The dependency chain at a glance}

The upper-bound and lower-bound branches are logically independent:
\[
\begin{array}{c}
\text{homogeneous GAS}
 \Longrightarrow \text{finite rational Lyapunov certificate}
 \Longrightarrow \HGAS_{\Q}(n,d)\text{ is c.e.}
\\[4pt]
\hline\\[-9pt]
\text{universal halting}
 \Longrightarrow \text{fixed DQL}
 \Longrightarrow \boxed{\text{modular/Jacobi source assembly}}
\\
\Longrightarrow \text{compact accumulator sign}
 \Longrightarrow \text{homogeneous wrapper}
 \Longrightarrow \text{cubicization and padding}.
\end{array}
\]
Everything outside the box is already represented by substantial proved
modules.  The box contains the remaining selector/classification and
canonical-source work, described precisely in Sections
\ref{sec:open-classification} and \ref{sec:open-source}.  The finite
unary--Jacobi code inside that future source is already compiled, including
both its fifty-dimensional static compact-pipeline instance and its
fifty-eight-dimensional dynamically controlled instances.  For the latter,
both a valid all-affine layout and the preferred geometry-faithful mixed
layout are compiled; the box denotes the missing dynamical assembly, not a
missing sparse-polynomial translation.

\section{The exact decision problem}
\label{sec:problem}

\subsection{Finite rational syntax}

For fixed dimension \(n\), a polynomial field code is a finite list of
rational monomials in each of \(n\) output coordinates.  Decoding gives a
map
\[
 F:\R^n\longrightarrow\R^n.
\]
The syntax layer is finite: rational numbers, exponent vectors, and lists.
The semantic layer contains real analysis, ODEs, and measures.  The
separation is deliberate.  A compiler may inspect only the finite syntax;
noncomputable objects such as flows or invariant measures may occur only in
correctness proofs.

The basic files are
\leanmodule{HGAS/Foundation/FiniteCode.lean} and
\leanmodule{HGAS/Foundation/FinitePolynomialCode.lean}.  The exact
fixed-resource language is defined in
\leanmodule{HGAS/Computability/HGASLanguage.lean} by
\[
\leanname{HGASQ n d F}
\quad:\Longleftrightarrow\quad
\begin{cases}
\text{every decoded coordinate of \(F\) is homogeneous of degree \(d\),}\\
0\text{ is globally asymptotically stable for the realized field.}
\end{cases}
\]
Malformed codes are negative instances.  The Boolean homogeneity test and
its correctness are
\leanname{PolynomialFieldCode.homogeneous} and
\leanname{PolynomialFieldCode.homogeneous_eq_true_iff}; its primitive
recursiveness is
\leanname{PolynomialFieldCode.homogeneous_primrec}.

\subsection{GAS without a hidden completeness assumption}

A polynomial ODE can have finite-time blow-up.  Thus one must not define
GAS using a global flow that may not exist.  The file
\leanmodule{HGAS/Dynamics/Stability.lean} starts instead with:
\begin{align*}
\leanname{IsForwardTrajectory}(F,\gamma,x_0)
 &:\Longleftrightarrow
 \gamma(0)=x_0\ \text{and}\ \dot\gamma(t)=F(\gamma(t))
 \text{ for }t\geq0,\\
\leanname{ForwardComplete}(F)
 &:\Longleftrightarrow
 \text{every initial point has a forward trajectory},\\
\leanname{LyapunovStable}(F,0)
 &:\Longleftrightarrow
 \text{nearby initial points stay nearby},\\
\leanname{GloballyAttractive}(F,0)
 &:\Longleftrightarrow
 \text{every forward trajectory tends to \(0\)}.
\end{align*}
The structure \leanname{GloballyAsymptoticallyStable F 0} contains all
three properties.  Consequently, every later GAS theorem has genuinely
proved forward completeness rather than assuming it through notation.

\subsection{Polar and angular equations}

Suppose
\[
F(r x)=r^{k+1}F(x),\qquad r\geq0.
\]
On the unit sphere define
\[
 a(u)=\langle u,F(u)\rangle,\qquad
 Y(u)=F(u)-a(u)u.
\]
Then \(\langle u,Y(u)\rangle=0\), so \(Y\) is tangent to the sphere.  If
\(x(t)=r(t)u(t)\), direct differentiation gives
\[
 \dot r=r^{k+1}a(u),\qquad
 \dot u=r^kY(u).
\]
With angular time \(s\), defined by \(\dd s/\dd t=r^k\), these become
\[
 \frac{\dd u}{\dd s}=Y(u),\qquad
 \frac{\dd}{\dd s}\log r=a(u).
\]
The second identity is the conceptual heart of the reduction.  Radial
contraction is controlled by additive orbit integrals of \(a\) along the
compact angular flow:
\[
\log\frac{r(s)}{r(0)}
   =\int_0^s a(\phi_\sigma u_0)\,\dd\sigma.
\]

The algebraic definitions
\leanname{innerRadialReward} and \leanname{innerAngularField} occur in
\leanmodule{HGAS/Dynamics/AngularGAS.lean}; the physical and angular polar
equations are proved there, including the logarithmic-radius derivative and
the conversion back to physical time.

\subsection{Uniform negative blocks and invariant probabilities}

On a compact flow \(\phi\) with continuous reward \(a\), say that there is
a uniform negative block if some \(T,c>0\) satisfy
\[
 \int_0^T a(\phi_t x)\,\dd t\leq-c
 \qquad\text{for every }x.
\]
If no such block exists, increasingly long almost-nonnegative orbit
segments produce occupation probabilities.  Compactness gives a weak
limit; the endpoint error vanishes after division by the segment length,
so the limit is invariant and has nonnegative \(a\)-average.  Conversely,
integrating a uniform block against an invariant probability gives a
strictly negative average.

This argument is fully formalized in
\leanmodule{HGAS/Dynamics/UniformNegativeBlocks.lean}.  Its key equivalences
are
\leanname{invariantAverage_neg_iff_uniformNegativeBlock} and
\leanname{maximalInvariantAverage_neg_iff_uniformNegativeBlock}.  Existence
of an invariant probability on a compact flow is proved by a
Krylov--Bogolyubov argument in
\leanmodule{HGAS/Dynamics/KrylovBogolyubov.lean}; compactness and attainment
of extrema appear in
\leanmodule{HGAS/Dynamics/InvariantProbabilityCompact.lean}.

\subsection{Ergodic maximizers for the full real flow}
\label{sec:ergodic-maximizer}

For the compact-sign argument later in the construction, it is not enough
merely to know that the largest invariant average is attained.  We want a
maximizer satisfying the zero--one law for invariant measurable sets.
The project uses the explicit full-flow definition
\[
\leanname{IsErgodicInvariantProbability}(\phi,\mu):
\quad
\begin{cases}
\mu\text{ is invariant under every }\phi_t,\\
\mu(A)\in\{0,1\}\text{ whenever }
  \phi_t^{-1}A=A\text{ modulo }\mu\text{ for every }t\in\R.
\end{cases}
\]
This definition and its zero--one consequences are in
\leanmodule{HGAS/Dynamics/ErgodicInvariantProbability.lean}.  In particular,
\leanname{IsErgodicInvariantProbability.zero_or_one_of_isInvariant} applies
directly to a measurable set that is literally invariant under a two-sided
flow.

There is a small but important pitfall here: applying a theorem about the
single map \(\phi_1\) would prove time-one ergodicity, not ergodicity for the
whole real action.  The two notions need not agree.  The module
\leanmodule{HGAS/Dynamics/ErgodicMaximizer.lean} therefore proves the
full-flow result directly.

The proof follows a convex-geometric route.  Send a probability measure to
its integration functional
\[
\mu\longmapsto
\left[f\longmapsto\int f\,\dd\mu\right]
\]
in the real weak dual of the bounded continuous functions.  On a compact
metric space this map is continuous, affine, and injective; injectivity is
the uniqueness part of real Riesz--Markov.  Since \Mathlib{} defines the
weak-star topology through evaluation coordinates but does not register the
needed locally convex instance in this form, the module proves that the
coordinate map is inducing and constructs that instance explicitly.  The
image of the invariant probabilities is then compact and convex.  For a
fixed observable \(f\), the measures maximizing
\(\int f\,\dd\mu\) form a nonempty compact face.  Krein--Milman supplies an
extreme point of that face, hence an extreme point of the full invariant
set.  The formal proof then:
\begin{enumerate}
\item converts real weak-dual extremality back to extremality in the positive
measure cone; and
\item proves directly that an extreme probability invariant under every
time map satisfies the full-flow zero--one law, by splitting it into
conditional probabilities on an invariant set and its complement.
\end{enumerate}

The resulting endpoint is
\leanname{exists_ergodicInvariantProbability_isMaxOn_average}: every bounded
continuous observable on a nonempty compact metric flow has an invariant
probability that is both ergodic for the whole real flow and maximizing
among \emph{all} invariant probabilities.  The chosen version is
\leanname{ergodicMaximizingInvariantProbabilityMeasure}.  Two useful
consequences are
\leanname{invariantAverage_le_of_forall_ergodic} and
\leanname{forall_invariantAverage_le_iff_forall_ergodic}: to prove an upper
bound for every invariant average, it is enough to prove it for ergodic
invariant probabilities.  This argument uses one ergodic maximizer and does
not require a measurable ergodic-decomposition theorem.

Finally,
\leanmodule{HGAS/Dynamics/ErgodicMaximalInvariantAverage.lean} identifies the
chosen ergodic average with the numerical maximum used elsewhere:
\[
\leanname{ergodicMaximizingInvariantProbabilityMeasure_average_eq},
\qquad
\leanname{maximalInvariantAverage_le_of_forall_ergodic}.
\]
These are unconditional kernel-checked theorems, not part of the remaining
horocycle classification interface.

Combining the angular ODE analysis and compact invariant-measure argument
gives the exact reusable criterion
\[
0\text{ is GAS for }F
\quad\Longleftrightarrow\quad
\max_{\mu\in\mathcal I(\phi)}
  \int a\,\dd\mu<0.
\]
The Lean theorem is
\leanname{globallyAsymptoticallyStable_iff_maximalAngularAverage_neg} in
\leanmodule{HGAS/Wrapper/AngularAverageGAS.lean}.  Its hypotheses say that
the finite-dimensional field is \(C^2\), homogeneous of positive degree,
and vanishes at the origin.

\section{The computably enumerable upper bound}
\label{sec:upper}

The upper-bound theorem is already unconditional:
\[
\boxed{\text{if \(d\) is odd, then \(\HGAS_{\Q}(n,d)\) is c.e.}}
\]
The exact Lean declaration is
\leanname{hgasQ_re_of_oddDegree} in
\leanmodule{HGAS/Lyapunov/RationalHomogeneousConverse.lean}.

\subsection{What a certificate contains}

Write
\[
\rho(x)=x_1^2+\cdots+x_n^2.
\]
A candidate stores a rational polynomial \(p\), a denominator exponent
\(r\), and exact positivity-margin data.  The associated degree-two
rational Lyapunov function is
\[
V(x)=\frac{p(x)}{\rho(x)^r}.
\]
The numerator of \(-\dot V\), after multiplying by the positive
denominator, is
\[
q(F,p,r)
=-\rho\,\nabla p\cdot F+2r\,p\,x\cdot F.
\]
This literal polynomial is
\leanname{rationalLyapunovNumerator} in
\leanmodule{HGAS/Lyapunov/RationalCertificateAlgebra.lean}.

The Boolean checker \leanname{checkRationalLyapunov} verifies two strict
polynomial inequalities by finite rational arithmetic: one for positivity
of \(p/\rho^r\), and one for strict decay.  The checker is primitive
recursive
(\leanname{checkRationalLyapunov_primrec}).  Its soundness theorem proves
not only stability and attraction but also forward completeness, using
coercive sublevel bounds; see
\leanmodule{HGAS/Lyapunov/CheckedRationalCertificateSoundness.lean}.

\subsection{Why every odd-degree GAS field has a certificate}

The completeness proof follows a geometric route.

\begin{enumerate}
\item The angular-flow construction produces a strict \(C^1\),
degree-two homogeneous Lyapunov function.  This is developed in
\leanmodule{HGAS/Lyapunov/HomogeneousConverseC1.lean} and
\leanmodule{HGAS/Lyapunov/EuclideanAngularConverse.lean}.

\item When \(d\) is odd, \(F(-x)=-F(x)\).  Averaging the Lyapunov function
with its antipodal pullback makes it even while preserving strict
positivity and strict decay.  The reusable implementation is in
\leanmodule{HGAS/Analysis/EvenC1Lyapunov.lean}.

\item On the unit sphere, approximate the function and its first
derivatives simultaneously by multivariate Bernstein polynomials.  The
development constructs the tensor Bernstein operator, proves derivative
formulas, and controls the \(C^1\) error in the modules
\leanmodule{HGAS/Analysis/TensorBernsteinPolynomial.lean},
\leanmodule{HGAS/Analysis/TensorBernsteinDerivative.lean}, and
\leanmodule{HGAS/Analysis/TensorBernsteinC1Approximation.lean}.

\item Evenness permits homogeneous polynomial reconstruction.  A sufficiently
small rational perturbation of the coefficients preserves the two strict
margins on the compact sphere.

\item Exact grid bounds replace semantic positivity with a finite
certificate.  The files
\leanmodule{HGAS/Lyapunov/ExactGridPositivity.lean} and
\leanmodule{HGAS/Lyapunov/GridRationalCertificateCompleteness.lean}
prove that a sufficiently fine rational grid, together with derivative
bounds, certifies positivity everywhere.
\end{enumerate}

The theorem
\leanname{realRationalLyapunovWitness_of_GAS_oddDegree} packages these
steps.  Then an enumeration of all finite candidates, combined with the
primitive-recursive checker, proves c.e. membership.

This proof is related to the homogeneous converse theorem of
Rosier~\cite{Rosier1992} and the rational Lyapunov results of Ahmadi and El
Khadir~\cite{AhmadiElKhadir2020}.  In the current Lean development, however,
the needed analytic and approximation chain is proved internally; the
literature theorem is not introduced as an axiom.

\section{The unconditional arithmetic source}
\label{sec:arithmetic}

\subsection{Universal halting and \(\Sigma^0_1\)-completeness}

The file \leanmodule{HGAS/Computability/SigmaOne.lean} defines
\[
\leanname{SigmaOneComplete}(P)
\quad:\Longleftrightarrow\quad
P\text{ is c.e. and every c.e. predicate many-one reduces to }P.
\]
It also defines a universal halting predicate for \Mathlib{} partial
recursive codes and proves
\leanname{universalHalts_sigmaOneComplete}.  The reduction from an arbitrary
c.e. set is explicit: choose a code for its partial characteristic
computation and pair that fixed code with the input.

\subsection{From primitive recursion to Diophantine equations}

The formal MRDP route is not a one-line appeal to the classical theorem.
It proves, by structural induction, that the graph of every unary
primitive-recursive function is Diophantine.  The difficult recursion case
uses a sparse radix encoding of an entire finite computation trace.

At a high level, a trace \(a_0,\ldots,a_N\) is packed into the digits of a
large integer.  A Diophantine certificate asserts:
\[
a_0=\mathrm{initial}(z),\qquad
a_{i+1}=\mathrm{step}(z,i,a_i)\ (i<N),\qquad
a_N=y.
\]
The bounded universal quantifier over \(i<N\) is itself compiled into a
sparse cipher certificate.  The common radix is chosen large enough that
all registers at all time positions are valid digits.

The main files are
\leanmodule{HGAS/Arithmetic/SparseBoundedForall.lean} and
\leanmodule{HGAS/Arithmetic/PrimitiveRecursiveDiophantine.lean}.  Principal
theorems include
\leanname{cipherCertificate_iff_boundedForall},
\leanname{boundedForall_dioph},
\leanname{recursionCertificate_iff},
\leanname{natPrimrec_unaryGraph_dioph}, and
\leanname{binaryPrimrecRel_dioph}.

The former interface
``every binary primitive-recursive relation is Diophantine'' is discharged
by
\leanname{UniversalHalting.binaryPrimitiveRecursiveRelationsAreDiophantine}
in \leanmodule{HGAS/Arithmetic/UnconditionalMRDP.lean}.  This yields a
fixed finite integer polynomial representing universal halting and the
semantic theorem
\leanname{unconditionalUniversalHaltingRepresentation_spec}.

The classical lineage is the Davis--Putnam--Robinson--Matiyasevich theorem
\cite{DavisPutnamRobinson1961,Matiyasevich1970,Matiyasevich1971,
Matiyasevich1993}.  The Lean proof makes its particular
primitive-recursive closure route explicit.  One scope point matters: the
final finite representative is chosen from a proof in \(\mathrm{Prop}\).
Thus its equation and variable counts are fixed natural numbers, but the
development does not print their numerical values or give a uniform
extractor for arbitrary c.e. programs.

\subsection{Diagonal--quadratic--linear form}

A DQL system has integer variables \(x_j\) and rows
\[
\sum_j\bigl(A_{ij}x_j^2+B_{ij}x_j\bigr)=c_i.
\]
No mixed quadratic term is allowed.  This form is expressive because a
multiplication gate \(xy=z\) can be replaced by
\[
w=x+y,\qquad w^2-x^2-y^2-2z=0,
\]
and natural-number variables can be represented over \(\Z\) by Lagrange's
four-square theorem~\cite{Lagrange1770}.  These two elementary equivalences are proved in
\leanmodule{HGAS/Arithmetic/DQL.lean}.

The finite circuit compiler preserves one fixed target arity after the
universal polynomial has been chosen.  In
\leanmodule{HGAS/Arithmetic/UnconditionalFixedDQL.lean}, the exact results
are:
\begin{align*}
&\leanname{unconditionalUniversalHaltingFixedDQL_solvable_iff},\\
&\leanname{unconditionalUniversalHaltingFixedDQL_computable},\\
&\leanname{unconditionalFixedDQLSolvability_sigmaOneComplete}.
\end{align*}
Thus solvability on one fixed DQL type is unconditionally
\(\Sigma^0_1\)-complete.

\section{Clocks, recurrence, and invariant measures}
\label{sec:clocks}

The reduction needs polynomial clocks whose long-term frequencies are
rigid under \emph{every} invariant probability, not merely under a selected
trajectory.

\subsection{Controller circles}

For a parameter \(p>1\), a two-dimensional controller has two hyperbolic
eigendirections.  Its normalized linear flow on the circle has only two
stationary directions; every other orbit moves monotonically through a
height coordinate.  Pairwise-disjoint translates of a strip then show that
an invariant probability gives the nonstationary set zero mass.

The explicit circle flow is \leanname{controllerFlow} in
\leanmodule{HGAS/Clocks/ControllerFlow.lean}.  The measure-theoretic endpoint
is
\leanname{invariantProbability_polynomialEquilibria_compl_eq_zero}; the
locked coordinate squares are identified by
\leanname{invariantProbability_lockedCoordinateSquares_compl_eq_zero}.
The underlying polynomial identities are in
\leanmodule{HGAS/Clocks/Controller.lean}.

The corresponding finite polynomial compiler is now also explicit.  For a
rational code \(p\), the state has four real coordinates
\((r,s,c,d)\), arranged as a controller circle and a driven phase circle:
\[
\begin{aligned}
\dot r&=-s(pr^2-1),&
\dot s&= r(pr^2-1),\\
\dot c&=-(pr)d,&
\dot d&=(pr)c.
\end{aligned}
\tag{5.1}
\]
The first pair is tangent to \(r^2+s^2=1\), and the second pair is tangent
to \(c^2+d^2=1\).  At a locked point \(pr^2=1\), the phase frequency is
\(pr=\pm\sqrt p\), which is the radical frequency used below.

The module \leanmodule{HGAS/Clocks/ControllerPolynomialCode.lean} fixes the
coordinate order, supplies inverse structured-state conversions, and emits
\leanname{ControllerPolynomialCode.fieldCode}.  Its theorem
\leanname{ControllerPolynomialCode.realize_fieldCode_arbitrary} proves
(5.1) at every point of \(\R^4\), without assuming either circle equation.
Theorems \leanname{ControllerPolynomialCode.controller_tangent} and
\leanname{ControllerPolynomialCode.phase_tangent} establish the two
tangency identities.  The emitted code is primitive recursive in \(p\) by
\leanname{ControllerPolynomialCode.fieldCode_primrec}, and
\leanname{ControllerPolynomialCode.fieldCode_termsDegreeLE} gives the
literal uniform raw-degree bound \(3\): the controller coordinates are
cubic, whereas the two phase coordinates are quadratic.

The finite compiler itself accepts every rational \(p\).  Positivity,
primality, and \(p>1\) enter only when the analytic support and radical
independence theorems are applied.

There is now also a genuine global flow for any finite family of these
coupled blocks.  The reusable construction in
\leanmodule{HGAS/Dynamics/ContinuousPhaseSkewFlow.lean} starts from a
continuous base flow \(\phi\) and continuous angular velocities \(a_i\).
It sets
\[
 A_i(t,x)=\int_0^t a_i(\phi_sx)\,\dd s,\qquad
 (x,z)\longmapsto
 \bigl(\phi_t x,(e^{\,iA_i(t,x)}z_i)_i\bigr).
\]
Joint continuity of the orbit integral and its cocycle identity give a
continuous real flow.  The theorem
\leanname{hasDerivAt_phaseSkewFlow_coe} proves the literal differential
equation
\[
 \dot z_i(t)=i\,a_i(\phi_t x)z_i(t)
\]
at the current base state.

The specialization
\leanmodule{HGAS/Clocks/ControllerPhaseFlow.lean} takes the product of the
explicit controller flows and \(a_i=p_i r_i\).  Its
\leanname{controllerPhaseFlow} therefore solves all four equations in
(5.1) simultaneously.  This is checked coordinatewise by
\leanname{controllerPhaseFlow_solves_controllerField} and
\leanname{controllerPhaseFlow_solves_lockedPhaseField}; each controller
projection is an exact continuous equivariant factor.  Thus the
time-dependent-looking phase equation is not merely a formal vector-field
expression: it is part of a global autonomous flow.

The analytic state orders all controller circles and then all phase
circles, whereas the finite two-controller compiler introduced in
Section~\ref{sec:dynamic-linked-compiler} stores two interleaved
\((r,s,c,d)\) blocks before the linked state.  The equations have been
proved in both representations, but the final source assembly must still
package the regrouping map and all downstream factor maps in one theorem.

\subsection{Ergodicity selects one controller equilibrium}

The support theorem above says that an invariant controller probability is
concentrated on a finite set of fixed points.  Ergodicity supplies the
missing selection principle.  The reusable theorem
\leanname{exists_eq_diracProba_of_isErgodicInvariantProbability_ae_fixed}
in \leanmodule{HGAS/Dynamics/ErgodicFixedSupport.lean} states:
\[
\begin{split}
 &\mu\text{ is an ergodic invariant probability, and}\\
 &\mu\text{-almost every point is fixed by every time map}
 \quad\Longrightarrow\quad
 \mu=\delta_x\text{ for some }x .
\end{split}
\tag{5.2}
\]
Indeed, modulo the null set of nonfixed points, every measurable set is
invariant.  Ergodicity makes \(\mu\) a zero--one measure, and the
standard-Borel zero--one theorem makes such a nonzero probability a Dirac
measure.

Applying (5.2) to the wandering-strip support theorem gives
\leanname{exists_eq_diracProba_at_controllerEquilibrium} in
\leanmodule{HGAS/Clocks/ErgodicControllerSelection.lean}.  It selects one
of the actual polynomial equilibria, not merely an unspecified fixed point.
The manuscript-facing consequences are
\[
 p\,r^2=1,\qquad r^2=\frac1p,\qquad
 s^2=1-\frac1p,\qquad (pr)^2=p .
\tag{5.3}
\]
The last identity says that the selected signed phase frequency is
\(\pm\sqrt p\).  The same module proves factor versions: if a continuous
equivariant observation of a larger ergodic flow is a controller circle,
then its pushed controller marginal is a Dirac mass at one equilibrium.
For a finite family, these componentwise statements can be intersected to
obtain one almost-sure controller tuple.  The next argument performs that
intersection and derives the phase law without assuming false global
equivariance.

\subsection{Radically independent torus rotations}

At a locked controller, the phase frequency is a signed square root of a
prime.  Distinct square roots of primes are linearly independent over
\(\Q\).  The formal proof builds the tower of fields generated by finitely
many prime square roots.  It proves by a square-class argument that a new
prime square root cannot lie in the preceding field, and then removes one
coefficient at a time from a putative rational relation.  Its final forms are
\leanname{linearIndependent_sqrt_primes} and
\leanname{signedPrimeSqrtFrequency_hasNoNontrivialIntegerCharacter} in
\leanmodule{HGAS/Clocks/RadicalIndependence.lean}.

For a torus translation, invariance annihilates every nontrivial Fourier
character.  Since characters determine probability measures, the only
invariant probability is Haar measure.  Theorems
\leanname{invariantProbability_eq_torusHaar_of_noCharacter} and
\leanname{signedPrimeSqrt_uniqueInvariantProbability} in
\leanmodule{HGAS/Clocks/RadicalClock.lean} formalize this argument.

One normalization is easy to miss.  The polynomial phase equation
\[
 \dot c=-\omega d,\qquad \dot d=\omega c
\]
uses angular velocity \(\omega\) in radians, while
\(\operatorname{UnitAddCircle}=\R/\Z\) measures turns.  Under the canonical
map \(x\mapsto e^{2\pi i x}\), the corresponding translation frequency is
\[
 v=\frac{\omega}{2\pi}.
\tag{5.4}
\]
The module \leanmodule{HGAS/Clocks/PhaseClockFlow.lean} proves this exact
semiconjugacy and separately proves that the physical rotation
\(z(t)=e^{i\omega t}z(0)\) solves the real polynomial phase field, with the
displayed sign.  The module
\leanmodule{HGAS/Clocks/ScaledRadicalClock.lean} proves that multiplying
every frequency by one common nonzero real scalar preserves the absence of
integer resonances.  Since \(1/(2\pi)\ne0\), division by \(2\pi\) does not
destroy radical independence.

The coordinatewise result in
\leanmodule{HGAS/Clocks/PhaseTorusFlow.lean} is
\leanname{physicalSignedRadicalPhase_equivariantMarginal_eq_haar}.  Given
an invariant ambient probability and an equivariant projection to physical
circle rotations with signed prime-square-root angular velocities, its
pushforward through the inverse angle map is normalized Haar measure on
the additive torus.  Ambient ergodicity is not required.  This theorem
fully discharges the radians-versus-turns issue for a phase projection that
is already known to rotate at fixed signed radical frequencies.

The missing passage from variable to fixed frequency is now proved in
\leanmodule{HGAS/Clocks/ErgodicControllerPhaseHaar.lean}.  For an ergodic
invariant probability \(\mu\) of the literal coupled flow,
\leanname{exists_lockedControllerTuple} selects a tuple \(x\) of controller
equilibria and proves simultaneously
\[
 (\pi_{\rm controller})_\ast\mu=\delta_x,
 \qquad
 z_{\rm controller}=x\quad\text{for \(\mu\)-almost every }z .
\tag{5.5}
\]
The finiteness of the clock index set is essential here: it lets the proof
intersect the full-measure coordinate fibres.

Phase projection is \emph{not} globally equivariant with a constant
rotation, because an unlocked controller has time-varying velocity
\(p_i r_i(t)\).  Instead,
\leanname{phaseProjection_ae_semiconj} proves, for every fixed time \(t\),
the almost-everywhere identity
\[
 \pi_{\rm phase}(\Phi_t z)
 =
 \operatorname{Rot}_{\,t(p_i\operatorname{Re}x_i)_i}
   (\pi_{\rm phase}z)
 \quad\text{for \(\mu\)-almost every }z .
\tag{5.6}
\]
The construction \leanname{lockedPhaseInvariantProbability} uses this
identity through equality of pushforward measures: it replaces the two
maps only almost everywhere, applies invariance of \(\mu\), and concludes
that the physical phase marginal is invariant under the constant locked
rotation.  No global semiconjugacy is smuggled into this step.

For distinct prime parameters, the locked frequencies are exactly
\(\sigma_i\sqrt{p_i}\), with \(\sigma_i\in\{+1,-1\}\).
The endpoint
\leanname{exists_lockedController_sign_and_phaseHaar} packages the selected
controller tuple, its joint Dirac marginal, almost-sure concentration on
the locked fibre, the sign tuple, the exact radical-frequency identity,
and Haar probability for the phase marginal after division by \(2\pi\).
Thus the coupled controller--phase invariant-measure classification itself
is complete.  A later source theorem must still exhibit the relevant
coupled clock flow as a factor of the full modular/Jacobi/detector
polynomial construction.

\subsection{Reusable measure transport}

Several supporting modules are worth separating from the application.

\begin{itemize}
\item \leanmodule{HGAS/Dynamics/InvariantProbability.lean} defines invariant
probabilities and equivariant pushforward.
\item \leanmodule{HGAS/Dynamics/PoincareRecurrence.lean} proves almost-everywhere
and existential forward recurrence for invariant probabilities.
\item \leanmodule{HGAS/Dynamics/RecurrenceTransport.lean} transports
recurrence through embeddings, homeomorphisms, and semiconjugacies.
\item \leanmodule{HGAS/Dynamics/FiniteSkewProductMeasure.lean} lifts an
invariant probability across a finite permutation skew product and
preserves the required recurrence statement.
\item \leanmodule{HGAS/Dynamics/InvariantMeasureTimeChange.lean} constructs
the normalized reciprocal-density measure for a positive generator time
change and proves the exact drift identity.
\end{itemize}

These constructions formalize the recurrence and invariant-measure tools
used in the manuscript; Walters~\cite{Walters1982} is a standard reference
for the surrounding ergodic theory.

\subsection{The detector theorem with physical clocks}

The arithmetic detector is naturally a function of additive turn
coordinates, whereas the polynomial source supplies physical complex
circles.  The theorem
\leanname{maximalInvariantAverage_neg_halfDetector_physicalClocks_iff_solvable}
in \leanmodule{HGAS/Arithmetic/DetectorPhysicalClockAverage.lean} composes
the inverse angle map (5.4), radical nonresonance, and Fourier
constant-term extraction.  Its hypotheses are a compact metric flow
\(\phi\), distinct prime parameters and signs, and a continuous equivariant
physical-phase projection.  Under precisely those hypotheses,
\[
 \mathcal M_\phi\!\left(
   -\operatorname{Re}\!
     \left(\text{expanded detector}_{\,1/2,\,1/2}\right)\right)<0
 \quad\Longleftrightarrow\quad
 \text{the given DQL system is solvable}.
\tag{5.7}
\]
There is no additional assumption saying that every phase marginal is
Haar: \leanname{physicalSignedRadicalPhase_equivariantMarginal_eq_haar}
proves it for every invariant joining.

The fixed-resource specialization
\leanname{maximalInvariantAverage_unconditionalUniversalHalting_physicalClocks_iff}
in
\leanmodule{HGAS/Arithmetic/UnconditionalFixedDetectorPhysicalClockAverage.lean}
replaces DQL solvability by the project's unary universal-halting
predicate.  The phase index type, prime family, source flow, and factor map
are fixed across inputs; only the finite detector coefficients vary.
This is already an unconditional arithmetic-to-average theorem, but it is
still a theorem \emph{about any flow carrying the stated physical factor}.
It does not construct that factor inside the final rational polynomial
source.

\subsection{The right-hand-side oscillator}
\label{sec:rhs-oscillator}

The reduction also needs to replace an input-dependent Laurent character by
an ordinary state variable.  Let \(\nu_i(x)\) be real polynomial frequency
codes and \(a_i\in\Q\), and put
\[
 W(x)=\sum_i a_i\nu_i(x).
\]
The module
\leanmodule{HGAS/Clocks/RHSOscillatorPolynomialCode.lean} adjoins two real
coordinates \(R=X+iY\) to an arbitrary \(n\)-dimensional polynomial field
and emits
\[
 R'=-4iW(x)R,\qquad
 X'=4W(x)Y,\qquad Y'=-4W(x)X.
\tag{5.8}
\]
The theorem \leanname{RHSOscillatorPolynomialCode.realize_fieldCode}
proves these equations at every point of \(\R^{n+2}\), with no unit-circle
or controller-lock hypothesis, and
\leanname{RHSOscillatorPolynomialCode.oscillator_tangent} proves
\(XX'+YY'=0\) identically.  The compiler is primitive recursive in the
base field, coefficients, and frequency codes.  If the base field has
degree at most \(D\) and every \(\nu_i\) has degree at most \(E\), the new
coordinates have degree at most \(E+1\); in particular, linear clock
frequencies produce quadratic oscillator equations.

There is a complementary, exact invariant-measure theorem in
\leanmodule{HGAS/Clocks/RHSOscillatorGraph.lean}.  Write the base as the
additive torus \((\R/\Z)^d\), let \(a\in\Z^d\), and drive one extra
additive circle at frequency
\[
 -\langle a,\omega\rangle
   =-\sum_i a_i\omega_i .
\]
Then
\[
 \Xi(x,R)=R+\sum_i a_i x_i
\tag{5.9}
\]
is conserved.  The triangular homeomorphism
\((x,R)\mapsto(x,\Xi(x,R))\) conjugates the joint flow to the base
translation times the identity flow.  If \(\omega\) has no nontrivial
integer resonance, the theorem
\leanname{ergodicInvariantProbability_graphClassification} says that every
ergodic invariant probability has Haar base marginal and equals Haar
pushed onto one unique graph
\[
 R=\xi-\sum_i a_i x_i .
\]
The selected graph has full measure.  The module also proves an
equivariant-factor form for the oscillator marginal of an arbitrary
ergodic ambient joining, and a nonergodic factor theorem for the Haar base
marginal.

For the fourth-power convention in (5.8), the integer character is
\(a_i=4c_i\).  The specialized frequency and invariant are proved
literally:
\[
 \Omega_R=-4\sum_i c_i\omega_i,\qquad
 \Xi=R+4\sum_i c_i x_i.
\tag{5.10}
\]
An arbitrary rational number does not define multiplication on
\(\R/\Z\).  The reusable structure
\leanname{FourIntegralRationalCoefficients} therefore records a rational
vector \(c\), an integer vector representing \(4c\), and the equality
between them.  The theorem
\leanname{rationalRhsErgodicInvariantProbability_graphClassification}
is the corresponding maximal correct rational-character statement.

These results close the general measure theory and the standalone
oscillator analysis.  What remains open is to identify, in one
kernel-checked source theorem, the physical polynomial oscillator (5.8)
with the normalized additive graph model (5.9)--(5.10), and to instantiate
all equivariance and finite-cover maps for the complete modular/Jacobi
polynomial state; see Section~\ref{sec:open-classification}.

\section{Theta constants and modular differential identities}
\label{sec:theta}

\subsection{Canonical theta constants}

For \(\tau\) in the upper half-plane, the project defines the classical
theta constants by absolutely convergent series:
\[
\vartheta_2(\tau),\qquad
\vartheta_3(\tau),\qquad
\vartheta_4(\tau),
\]
and \(H_j=\vartheta_j^4\).  The definitions and comparison with
\Mathlib{}'s Jacobi theta series are in
\leanmodule{HGAS/Modular/ThetaConstants.lean}; holomorphy is proved in
\leanmodule{HGAS/Modular/ThetaHolomorphy.lean}.

The modular transformation and cusp-asymptotic modules support proofs by a
standard pattern: build a holomorphic modular expression, prove its
transformation under the generators \(S,T\), show that it vanishes at the
cusp, and use a low-weight vanishing theorem.  In this way the development
proves:
\begin{itemize}
\item Jacobi's quartic identity
\[
\vartheta_3^4=\vartheta_2^4+\vartheta_4^4
\]
as \leanname{jacobi_quartic_identity};
\item the theta derivative identities
\leanname{normalizedDeriv_H₂},
\leanname{normalizedDeriv_H₃}, and
\leanname{normalizedDeriv_H₄};
\item Ramanujan's equations
\leanname{ramanujan_E2}, \leanname{ramanujan_E4}, and
\leanname{ramanujan_E6}.
\end{itemize}
The corresponding modules are
\leanmodule{HGAS/Modular/ThetaJacobiIdentity.lean},
\leanmodule{HGAS/Modular/ThetaDerivativeIdentities.lean}, and
\leanmodule{HGAS/Modular/RamanujanIdentities.lean}.

These formulas are classical; useful references are
Chandrasekharan~\cite{Chandrasekharan1985}, Mumford~\cite{Mumford1983}, and
the DLMF~\cite{DLMF}.  The Lean files derive the normalizations they use
rather than importing the formulas as axioms.

\section{The Jacobi theta state}
\label{sec:jacobi}

\subsection{Residual quadrics and first-order Wronskians}

Fix \(\tau\in\mathbb H\), and write
\[
T(z)=\bigl(\vartheta_1(z|\tau),\vartheta_2(z|\tau),
           \vartheta_3(z|\tau),\vartheta_4(z|\tau)\bigr).
\]
Two homogeneous quadrics cut out the intended residual curve in
\(\C^4\).  The actual theta vector satisfies both; this is proved in
\leanmodule{HGAS/Modular/JacobiThetaResiduals.lean}.  The duplication
identities are formalized in
\leanmodule{HGAS/Modular/JacobiThetaDuplication.lean}.

The first-order closed elliptic system needs three Wronskian identities.
The proof in
\leanmodule{HGAS/Modular/JacobiThetaWronskians.lean} proceeds from the
series:
\begin{enumerate}
\item rewrite half-characteristic theta functions as summable series;
\item pair even and odd lattice indices;
\item prove doubling recurrences for the Wronskian and for the product of
theta constants;
\item identify the multiplicative constant from the cusp asymptotic;
\item use analyticity and cancellation to extend across zeros.
\end{enumerate}
The endpoint is \leanname{canonicalJacobiThetaWronskiansAt}.  No local
rank or nonvanishing hypothesis remains in that theorem.

\subsection{The Delta identity}
\label{sec:delta}

The canonical algebraic Jacobi state stores
\[
p_j=\vartheta_j(\tau)^2,\quad
H=\frac{E_2(\tau)}{12},\quad
T_j=\vartheta_j(z|\tau),\quad
\Delta=\frac{1}{\pi}\partial_z\vartheta_1(z|\tau),\quad
U=T_1^{-1}.
\]
All first-order coordinate equations except the \(\Delta\)-equation follow
from the Wronskians.  After clearing the reciprocal \(U\), the missing
equation is exactly
\[
T_1\,\partial_z^2T_1
=(\partial_zT_1)^2
 -\pi^2p_3p_4T_2^2
 -4\pi^2K\,T_1^2,
\qquad
K=H+\frac{p_3^2+p_4^2}{12}.
\tag{7.1}
\]
The signs in (7.1) are those used by the actual Lean vector field.

The module
\leanmodule{HGAS/Modular/JacobiThetaDeltaIdentity.lean} now proves this
identity unconditionally.  The proof is instructive.

\begin{enumerate}
\item Define the cross-multiplied error \(E_1\) for \(T_1\) and a companion
error \(E_2\) for \(T_2\).
\item Differentiate the first Wronskian.  The three Wronskians and the two
residual quadrics yield the polynomial identity
\[
T_1^2E_2=T_2^2E_1.
\]
\item Since \(T_1,T_2\) have no common zero, patch \(E_1/T_1^2\) and
\(E_2/T_2^2\) to an entire function of \(z\).
\item Differentiate theta quasiperiodicity.  The patched function has
periods \(1\) and \(\tau\), hence is bounded on a fundamental
parallelogram.  Liouville's theorem makes it constant.
\item Prove the heat equation
\[
\partial_z^2\vartheta_2
  =4\pi i\,\partial_\tau\vartheta_2
\]
by termwise differentiation of the Gaussian series, with summability and
local uniformity made explicit.
\item At \(z=0\), evenness gives \(\partial_z\vartheta_2=0\).  The heat
equation, the modular derivative identity, and Jacobi's quartic identity
show that the companion error is zero there.
\item The patched constant is therefore zero.  The patching identity
extends the conclusion across every zero of \(T_1\).
\end{enumerate}

The exported theorems are
\leanname{jacobiThetaDeltaIdentity},
\leanname{canonicalJacobiThetaState_elliptic_unconditional}, and
\leanname{canonicalJacobiThetaRealAffine_isScaledTrajectory_unconditional}.
The last theorem checks the normalization: restricting along
\(z(t)=z_0+tv\) gives elliptic speed \(v\pi\).

Brezhnev's closed elliptic system provides the literature context and
normalization comparison~\cite[Theorem 4.1 and equation (34),
pp.~702--703]{Brezhnev2013}.  The Lean proof of (7.1) is independent of
treating that reference as a formal assumption.

\subsection{The effective ten-coordinate Jacobi compiler}
\label{sec:jacobi-compiler}

The analytic theorems above describe complex-valued functions.  A
computability reduction, however, needs a finite list of rational
polynomials.  The three modules
\[
\begin{gathered}
\leanmodule{HGAS/Modular/JacobiPolynomialCode.lean},\\
\leanmodule{HGAS/Modular/JacobiPolynomialComputability.lean},\qquad
\leanmodule{HGAS/Modular/JacobiPolynomialDegrees.lean}
\end{gathered}
\]
now provide that finite compiler for the entire reciprocal Jacobi chart.

The complex coordinate order is fixed once and for all:
\[
\begin{array}{c|cccccccccc}
\text{index}&0&1&2&3&4&5&6&7&8&9\\ \hline
\text{coordinate}&p_2&p_3&p_4&H&T_1&T_2&T_3&T_4&\Delta&U .
\end{array}
\]
The maps \leanname{JacobiPolynomialCode.point} and
\leanname{JacobiPolynomialCode.state} move between this vector and the
structured \leanname{JacobiState}; \leanname{state_point} and
\leanname{point_state} prove that the two maps are inverse.  Thus a semantic
theorem about coordinate \(i\) cannot hide a permutation of the state.

The reusable syntax \leanname{ComplexPolynomialCode} represents one complex
polynomial by a pair of rational sparse polynomials in the real and
imaginary parts of all variables.  Complex addition and multiplication are
finite list operations.  A \leanname{GaussianRationalCode} is just the pair
of rational real and imaginary parts of a scalar.  Consequently no complex
floating-point arithmetic enters the compiler.

There is also a reusable realification operation for a family whose number
of equations need not equal its number of variables.  If
\[
H:\operatorname{Fin}(m)\longrightarrow
  \leanname{ComplexPolynomialCode}(n),
\]
then \leanname{ComplexPolynomialFamilyCode.realify} emits \(2m\) real
polynomials in \(2n\) real variables: first all real parts, then all
imaginary parts.  This genuinely extends field realification, for which the
number of outputs is constrained to equal \(n\).  The theorem
\leanname{ComplexPolynomialFamilyCode.complexifyPoint_eval_realify}
reassembles the \(2m\) evaluated scalars into the original \(m\) complex
values, and
\leanname{ComplexPolynomialFamilyCode.forall_eval_realify_eq_zero_iff}
proves
\[
\bigl[\text{all \(2m\) real scalar equations vanish}\bigr]
\quad\Longleftrightarrow\quad
\bigl[\text{all \(m\) complex equations vanish}\bigr].
\tag{7.2}
\]
The construction preserves every uniform raw-degree bound and is primitive
recursive, by
\leanname{ComplexPolynomialFamilyCode.realify_termsDegreeLE} and
\leanname{ComplexPolynomialFamilyCode.realify_primrec}.  These results live
in \leanmodule{HGAS/Algebra/ComplexPolynomialFamilyCode.lean}; they are not
specific to theta functions and can be reused for any finite complex
constraint family.

Two ten-component fields are emitted.  The first is the elliptic
derivation \(\mathcal Z\):
\[
\begin{aligned}
\mathcal Z(p_2,p_3,p_4,H)&=0,&
\mathcal Z(T_1)&=\Delta,\\
\mathcal Z(T_2)&=U(\Delta T_2-p_2T_3T_4),&
\mathcal Z(T_3)&=U(\Delta T_3-p_3T_4T_2),\\
\mathcal Z(T_4)&=U(\Delta T_4-p_4T_2T_3),&
\mathcal Z(U)&=-U^2\Delta,\\
\mathcal Z(\Delta)&=
 U(\Delta^2-p_3p_4T_2^2)-4KT_1,&
K&=H+\frac{p_3^2+p_4^2}{12}.
\end{aligned}
\]
The second is the modular derivation \(\Phi\).  Its coefficient block begins
\[
\begin{aligned}
\Phi(p_2)&=2\left(H+\frac{p_3^2+p_4^2}{12}\right)p_2,\\
\Phi(p_3)&=2\left(H+\frac{p_2^2-p_4^2}{12}\right)p_3,\\
\Phi(p_4)&=2\left(H-\frac{p_2^2+p_3^2}{12}\right)p_4,\\
\Phi(H)&=2H^2-\frac{p_2^4+p_3^4+p_4^4}{144}.
\end{aligned}
\]
The compiler also includes, without simplification, the complete
\(T_1,T_2,T_3,T_4,\Delta,U\) equations from
\leanname{JacobiState.modularVectorField}.  In particular, the
degree-eight term
\[
-\frac12 p_2p_3p_4T_2T_3T_4U^2
\]
is present in the \(\Delta\)-coordinate.  Retaining unsimplified terms is
intentional: later degree and homogeneity compilers reason about the literal
raw sparse lists, so their correctness cannot depend on a noncomputable
normal form.

The semantic checks are coordinatewise and exact:
\[
\begin{aligned}
\leanname{eval_ellipticFieldCode}&:
 \operatorname{eval}(\mathcal Z_{\rm code},z)=\mathcal Z(z),\\
\leanname{eval_modularFieldCode}&:
 \operatorname{eval}(\Phi_{\rm code},z)=\Phi(z).
\end{aligned}
\]
Given Gaussian-rational clocks \(\kappa_{\mathcal Z},\kappa_\Phi\),
\leanname{combinedFieldCode} forms
\(\kappa_{\mathcal Z}\mathcal Z+\kappa_\Phi\Phi\), and
\leanname{eval_combinedFieldCode} proves the corresponding structured-state
identity.

Realification then turns the ten complex coordinates into twenty real
coordinates.  The theorem
\leanname{complexify_realize_compileFieldCode_arbitrary} is stronger than a
chart-only statement: it holds at \emph{every} point of \(\R^{20}\).  After
reassembling that point as a complex ten-vector and then a
\leanname{JacobiState}, evaluation of the emitted real field is exactly the
combined complex field.  None of the residual equations is assumed for this
global semantic theorem.

The same file emits finite codes for the intrinsic chart residuals
\[
\begin{aligned}
R_0&=p_3^2-p_2^2-p_4^2,\\
R_1&=p_2T_4^2-p_4T_2^2-p_3T_1^2,\\
R_2&=p_2T_3^2-p_3T_2^2-p_4T_1^2,\\
R_U&=UT_1-1.
\end{aligned}
\]
Their order and semantics are fixed by \leanname{residualCode} and
\leanname{eval_residualCode}.

All code-producing maps are effective.  The fixed elliptic and modular
fields are primitive-recursive constants;
\leanname{combinedFieldCode_primrec} is primitive recursive in both clock
codes; \leanname{compileFieldCode_primrec} composes this construction with
primitive-recursive realification; and \leanname{residualCode_primrec}
handles the fixed four-element residual family.  Explicit computability
corollaries are exported for each.

Finally, the degree ledger is syntactic and therefore robust under duplicate
or cancelling terms:
\[
\begin{array}{c|c}
\text{finite family}&\text{raw degree bound}\\ \hline
\text{elliptic field}&5\\
\text{modular field}&8\\
\text{every two-clock combination and its realification}&8\\
\text{four residuals}&3.
\end{array}
\]
The exact theorem names are
\leanname{ellipticFieldCode_termsDegreeLE},
\leanname{modularFieldCode_termsDegreeLE},
\leanname{combinedFieldCode_termsDegreeLE},
\leanname{compileFieldCode_termsDegreeLE}, and
\leanname{residualCode_termsDegreeLE}.  The fixed ledgers are decided by
an executable raw-list predicate; the parameterized clock and realification
bounds are then propagated by reusable degree lemmas.  The constructive
decision and closure rules are reusable outside the Jacobi compiler; they
live in \leanmodule{HGAS/Accumulator/PolynomialDegrees.lean} and
\leanmodule{HGAS/Algebra/ComplexPolynomialDegrees.lean}.

This compiler closes a genuine effectivity gap: once the semantic selector
produces the appropriate Jacobi state and rational clock data, the complete
polynomial field is already finite, rational, degree-bounded, and
computable.  It does \emph{not} prove that every recurrent polynomial state
comes from the analytic theta construction, nor does it construct the
invariant-measure selector.  Those remain the semantic tasks in
Section~\ref{sec:open-classification}.

\subsection{The linked unary--Jacobi compiler}
\label{sec:linked-jacobi-compiler}

The reduction uses both a four-coordinate unary modular block
\((A,C,E,Y)\) and the ten-coordinate Jacobi block above.  The unary block
is convenient for modular classification, while the Jacobi block is
convenient for the polynomial theta selector.  They cannot merely be
identified informally: the compiler keeps them as independent variables
and imposes explicit polynomial links.

The module \leanmodule{HGAS/Modular/LinkedJacobiPolynomialCode.lean}
constructs the combined field on
\[
4+10=14\quad\text{complex coordinates},\qquad
2(4+10)=28\quad\text{real coordinates}.
\]
It uses the generic block-field constructors rather than copying any
coordinate formula.  The elliptic clock acts only on the Jacobi block; the
modular clock acts on both blocks.  Thus the only runtime field parameters
are two Gaussian-rational codes
\(\kappa_{\mathcal Z}\) and \(\kappa_\Phi\).  The theorem
\leanname{LinkedJacobiPolynomialCode.eval_complexFieldCode} proves exact
structured semantics, and the stronger theorem
\leanname{LinkedJacobiPolynomialCode.complexify_realize_compileFieldCode_arbitrary}
proves the real compiler correct at every \(x\in\R^{28}\), before any
constraint equation is assumed.

Ten complex residuals define the intended linked chart.  Four are the
intrinsic Jacobi residuals \(R_0,R_1,R_2,R_U\) listed above.  The other six
are
\[
\begin{gathered}
Y^4-A-C,\qquad p_2^2-A,\qquad p_3^2-Y^4,\qquad
p_4^2-C,\\
12H-E,\qquad p_3-Y^2.
\end{gathered}
\tag{7.3}
\]
The final equation chooses the signed square-root branch rather than
retaining only its square.  The executable enumeration
\leanname{LinkedJacobiPolynomialCode.residualIndex} is proved bijective
onto the ten semantic names, and
\leanname{LinkedJacobiPolynomialCode.eval_residualCode} proves that each
finite code evaluates to its named residual.

Applying the arbitrary-family construction (7.2) produces twenty rational
real residual codes in twenty-eight variables.  The theorem
\leanname{LinkedJacobiPolynomialCode.forall_eval_realResidualCode_eq_zero_iff}
says that all twenty scalar equations vanish exactly when every one of the
ten named complex residuals vanishes.  This is an equivalence, not merely
the easy implication from complex vanishing to real vanishing.

The entire linked compiler is effective:
\leanname{LinkedJacobiPolynomialCode.complexFieldCode_primrec} and
\leanname{LinkedJacobiPolynomialCode.compileFieldCode_primrec} are
primitive recursive in the two clock codes, while the ten complex and
twenty real residual families are fixed primitive-recursive finite
objects.  The syntactic degree ledger is
\[
\deg(\text{linked field})\leq8,\qquad
\deg(\text{each linked residual})\leq4,
\tag{7.4}
\]
both before and after realification.  The exact endpoints are
\leanname{LinkedJacobiPolynomialCode.compileFieldCode_termsDegreeLE} and
\leanname{LinkedJacobiPolynomialCode.realResidualCode_termsDegreeLE}.

This is a completed finite compiler for the linked block.  It does not say
that an arbitrary recurrent point satisfying the ten residuals is selected
by analytic theta data, and it does not dynamically couple the linked block
to arithmetic detectors.  The controller coupling is supplied by the next
module; the analytic selector remains a semantic source-assembly question,
not an omission from the sparse-polynomial compiler.

\subsection{The dynamic two-controller linked-Jacobi compiler}
\label{sec:dynamic-linked-compiler}

The module
\leanmodule{HGAS/Modular/TwoControllerLinkedJacobiPolynomialCode.lean}
now adjoins two complete four-real-coordinate controller--phase blocks to
the twenty-eight-real-coordinate linked state:
\[
 4+4+28=36.
\tag{7.5}
\]
The first controller supplies the modular angular velocity and the second
supplies the integer-like parameter.  At a structured state define
\[
 \omega=p_\Omega r_\Omega,\qquad
 n_{\rm dyn}=p_Nr_N,\qquad
 \kappa_{\mathcal Z}=\frac{\omega+n_{\rm dyn}}2,\qquad
 \kappa_\Phi=i\omega .
\tag{7.6}
\]
The linked compiler originally takes constant Gaussian-rational
coefficients.  To obtain (7.6) without pretending that a state-dependent
quantity is a coefficient, the new compiler forms two fixed basis fields,
with parameters \((1,0)\) and \((0,i)\), and multiplies their coordinates
by the real polynomial codes for \(\kappa_{\mathcal Z}\) and \(\omega\).
The coordinatewise identity
\leanname{point_combinedVectorField_dynamic_basis} proves that the result
is exactly the intended linked vector field with parameters (7.6).

The maps \leanname{TwoControllerLinkedJacobiPolynomialCode.point} and
\leanname{TwoControllerLinkedJacobiPolynomialCode.state} are inverse.
Consequently,
\leanname{TwoControllerLinkedJacobiPolynomialCode.realize_fieldCode_arbitrary}
proves the complete field semantics at every point of \(\R^{36}\), not
only on the two controller tori or on the linked residual variety.
The code-producing map is primitive recursive in the two rational
parameters \(p_\Omega,p_N\).

The degree bookkeeping distinguishes two notions.  In the controller
variables and the linked target variables separately, the complete field
has bidegree at most
\[
 (\deg_{\rm controller},\deg_{\rm linked})\leq(3,8).
\tag{7.7}
\]
The dynamically scaled linked coordinates themselves have bidegree at
most \((1,8)\).  Ordinary total degree is therefore at most \(9\), and
\emph{not} \(8\): multiplying a linked monomial of degree \(8\) by a
controller-linear clock genuinely produces degree \(9\).  These raw sparse
list bounds are theorems in
\leanmodule{HGAS/Modular/TwoControllerLinkedJacobiPolynomialDegrees.lean};
the reusable bidegree closure calculus is in
\leanmodule{HGAS/Algebra/PolynomialBlockDegrees.lean}.

This closes the dynamic finite-code coupling for the two clock parameters.
It still does not prove that the analytic recurrent selector lands on the
linked residual set, nor does it construct the physical detector and
right-hand-side-oscillator factor maps from the resulting compact source.

\subsection{Global projective saturation}
\label{sec:saturation}

The map
\[
\C\longrightarrow\mathbb{CP}^1,\qquad
z\longmapsto[\vartheta_1(z|\tau):\vartheta_2(z|\tau)]
\]
is globally onto.  Rather than rely on a derivative being nonzero, the
formal proof uses the complex open mapping theorem.

First, the two theta functions have no common zero and are not globally
proportional.  Their ratio is therefore open in either projective chart.
Second, level-two quasiperiodicity reduces every value to a compact
fundamental rectangle, so the range is compact and hence closed.  A
nonempty subset of \(\mathbb{CP}^1\) that is both open and closed is the
whole connected projective line.

The topology is isolated in the reusable module
\leanmodule{HGAS/Topology/ComplexProjectiveLine.lean}.  The theta-specific
result is then strengthened algebraically.  Given a nonzero vector on the
two residual quadrics, match its first two coordinates projectively.
Subtracting the residual equations forces each of the last two coordinates
to agree up to sign.  The four sign patterns are realized by \(z\),
\(\tau-z\), \(z+1\), and \((\tau-z)+1\), with the appropriate common
multiplier.

The exact endpoints in
\leanmodule{HGAS/Modular/JacobiThetaProjectiveSurjectivity.lean} are
\leanname{exists_smul_jacobiThetaVector_eq_of_residual},
\leanname{projectiveJacobiTheta_surjective_onto_residualCurve}, and
\leanname{range_projectiveJacobiTheta_eq_residualCurve}.  They are packaged
for the selector by
\leanname{canonicalJacobiProjectiveCoverage} and
\leanname{pointwiseJacobiBlochSaturation} in
\leanmodule{HGAS/Modular/JacobiPointwiseSaturationUnconditional.lean}.

This is a set-theoretic global saturation theorem.  It does not by itself
classify recurrent polynomial states, choose a measurable theta parameter,
or prove a statement about invariant measures.

The analytic ingredients in this open-and-closed argument are standard
one-variable complex analysis; see, for example,
Ahlfors~\cite[pp.~271--274]{Ahlfors1966}.

\section{The modular quotient and the remaining classification}
\label{sec:open-classification}

\subsection{What is already constructed}

The quotient
\[
\Gamma(2)\backslash\mathrm{SL}_2(\R)
\]
and its right-unipotent horocycle flow are constructed in
\leanmodule{HGAS/Modular/GammaTwoFrameQuotient.lean}.  The proof establishes
proper discontinuity of the arithmetic action, the quotient-covering
topology, continuity of right multiplication, every fixed nonzero time
scale, and projection to the upper-half-plane quotient.

For an arbitrary continuous real flow, the reusable module
\leanmodule{HGAS/Dynamics/PeriodicOrbitMeasure.lean} constructs normalized
arclength probability on a periodic orbit, proves invariance, proves the
integral formula, and proves uniqueness among invariant probabilities
carried by that orbit.

The file
\leanmodule{HGAS/Modular/HorocycleMeasureClassification.lean} defines the
precise alternatives:
\[
\text{ergodic invariant probability}
\quad=\quad
\begin{cases}
\text{normalized probability on a closed horocycle},\\
\text{right-invariant frame probability}.
\end{cases}
\]
All changes by a nonzero constant speed and all normalized periodic-orbit
conclusions are proved.

\subsection{The conditional recurrent-measure bridge}

The formal development now carries the pure classification hypothesis all
the way through an arbitrary state-space factor map.  For a state flow
\(\Psi\), a real scale \(s\), and an observation
\[
\pi:S\longrightarrow\Gamma(2)\backslash\mathrm{SL}_2(\R),
\]
the proposition \leanname{IsModularHorocycleFactorMap} records two concrete
facts: \(\pi\) is continuous, and
\[
\pi(\Psi_t x)=h_{st}(\pi(x))
\quad\text{for every \(t\in\R\) and \(x\in S\)}.
\tag{8.1}
\]
Here \(h\) denotes the unit-speed quotient horocycle flow.  The
\leanname{modularHorocycleMarginal} of a state probability is simply its
pushforward by \(\pi\).  Equivariance proves that the marginal of an
ergodic \(\Psi\)-invariant probability is an ergodic
\(h_{s\cdot}\)-invariant probability; this is
\leanname{IsModularHorocycleFactorMap.ergodicMarginal}.

Now assume, explicitly,
\[
h_{\rm DS}:\leanname{GammaTwoDaniSmillieTheorem},
\qquad s\ne0.
\]
Then
\leanname{GammaTwoDaniSmillieTheorem.classifyErgodicStateMarginal} gives a
complete dichotomy for the pushforward.  Either it is canonical normalized
arclength on a periodic horocycle, or it is the classified normalized
right-invariant frame probability.  In the first branch the conclusion is
stronger than a statement about the marginal: the original state measure
has mass one on the inverse image of that periodic orbit.  This support
property is packaged as \leanname{IsCarriedOverClosedHorocycle}.

There is also a recurrent form.  On a second-countable pseudometric Borel
state space,
\leanname{GammaTwoDaniSmillieTheorem.classifyErgodicRecurrentStateMarginal}
combines the dichotomy with Poincar\'e recurrence: for almost every state,
both its state orbit and its observed modular orbit are forward recurrent
at time zero.  Recurrence downstairs follows from (8.1) and continuity,
not from a separate recurrence assumption.

For globally trivial finite-sheet models
\[
\bigl(\Gamma(2)\backslash\mathrm{SL}_2(\R)\bigr)\times\operatorname{Fin}(n)
\]
with a finite permutation skew-product flow, the first projection is proved
to be such a factor map by
\leanname{finitePermutationSkewProductModularFactor}.  Consequently the
classification and recurrent classification specialize to theorems
\leanname{GammaTwoDaniSmillieTheorem.classifyErgodicFinitePermutationCover}
and
\leanname{GammaTwoDaniSmillieTheorem.classifyErgodicRecurrentFinitePermutationCover}.
Conversely,
\leanname{GammaTwoDaniSmillieTheorem.classifyAndLiftErgodicModularProbability}
lifts any ergodic modular probability to a chosen global sheet, preserves
its first marginal and ergodicity, proves almost-everywhere recurrence
upstairs, and applies the same conditional dichotomy.

Every implication in this subsection is kernel checked in
\leanmodule{HGAS/Modular/RecurrentMeasureClassification.lean}.  Its
hypotheses are essential and remain visible.  The file neither proves
\leanname{GammaTwoDaniSmillieTheorem} nor constructs the paper-specific
theta-state map \(\pi\) and proves (8.1).  Moreover, its finite-cover
specialization is the honest global product with \(\operatorname{Fin}(n)\);
it is not a measurable lifting theorem for an arbitrary nontrivial
connected finite cover.  Thus these are \statusinterface{} results: the
measure-theoretic transport is finished once the deep classification and
the concrete factor map are supplied.

\subsection{Open theorem A: Dani--Smillie/Hubbard--Miller}

\statusopen\quad The proposition
\leanname{GammaTwoDaniSmillieTheorem} has not yet been proved.  Closing it
requires finite covolume for \(\Gamma(2)\) and the invariant-measure
classification (or a strong enough equidistribution theorem) for the
horocycle flow.  This is the deep literature input of
Dani--Smillie~\cite{DaniSmillie1984} and
Hubbard--Miller~\cite{HubbardMiller2010}.

The project has also formalized the analytic core of the Haar-exclusion
argument.  In
\leanmodule{HGAS/Dynamics/CuspMeasureDivergence.lean}, a cusp measure with
density \(y^{-2}\dd y\) is shown to have divergent \(y^M\)-moment for
\(M\geq2\).  In \leanmodule{HGAS/Modular/HaarExclusion.lean}, theta cusp
asymptotics give the required lower bound, and arbitrary joinings are
handled by pushforward monotonicity of lower integrals.  The endpoint
\leanname{haarMarginal_reciprocalSlowdown_lintegral_eq_top} is factored by
an explicit cusp-marginal domination hypothesis.  Still open are the
Iwasawa cusp formula for the actual quotient Haar probability and the
wiring of the modular coordinate to the mixed compactification slowdown.

\subsection{Open theorem B: recurrent state-map and selector wiring}

\statusopen\quad Even after the pure horocycle theorem, the following
paper-specific bridges remain:
\begin{enumerate}
\item reconstruct the recurrent generic unary modular polynomial state
from a frame-quotient point, including multiplier and deck data;
\item prove properness/finite-cover properties of that concrete state map
and identify it with either the globally trivial finite-skew-product
interface above or a separately proved nontrivial-cover lifting theorem;
\item combine the unconditional projective Jacobi saturation with global
continuation of the commuting modular and elliptic derivations;
\item exclude nonreal Bloch escape and degenerate branches for recurrent
states;
\item connect the finite moment selector to invariant/time averages and
the reflected branch.
\end{enumerate}

Several pieces are already available.  Degenerate unary divisors are
handled in \leanmodule{HGAS/Modular/UnaryDegenerateRecurrence.lean}.
Cesaro extraction for absolutely summable exponential series is proved as
\leanname{naturalCesaroMean_exponentialSeries_tendsto}.  The exact
quadratic resonance table is in
\leanmodule{HGAS/Modular/JacobiMomentResonance.lean}, and the reflected
trace reindexing is
\leanname{reflectedJacobiTrace_eq_canonical}.  These theorems begin from an
explicit Bloch expansion; obtaining that expansion from an arbitrary
recurrent polynomial state is the missing selector theorem.

The relevant classical differential systems are discussed in
Brezhnev--Lyakhovich--Sharapov~\cite{BrezhnevLyakhovichSharapov2011} and
Brezhnev~\cite{Brezhnev2013}.  Their use in the final proof cannot be
replaced by a citation inside Lean: the required consequences must be
proved in the project, or derived from independently formalized versions.

\section{Compactification and the accumulator}
\label{sec:compact}

\subsection{Stereographic compactification}

For \(z\in\R^m\), set
\[
u(z)=\frac{2z}{1+\|z\|^2},\qquad
w(z)=\frac{1-\|z\|^2}{1+\|z\|^2}.
\]
Then \(\|u\|^2+w^2=1\), and \(z=u/(1+w)\) away from the point
\((u,w)=(0,-1)\).  If the affine drift is \(g\), differentiation gives the
polynomial tangent field
\[
\dot u=(1+w)g-(u\cdot g)u,\qquad
\dot w=-(1+w)(u\cdot g).
\]
The point at infinity is an equilibrium.  These identities are proved in
\leanmodule{HGAS/Compactification/Stereographic.lean}, principally
\leanname{stereo_mem_sphere}, \leanname{accumulator_tangent},
\leanname{stereoUDeriv_eq_field}, and
\leanname{stereoWDeriv_eq_field}.

\subsection{Mixed central clearing}

The source contains both affine blocks and blocks that are already compact.
Each affine block gains a central sphere coordinate \(w_b\).  A monomial is
homogenized blockwise, and the common slowdown
\[
\kappa\prod_b w_b^M
\]
clears every denominator.  The finite syntax and exact chart evaluation are
implemented in
\leanmodule{HGAS/Compactification/MixedCentralPolynomialCode.lean}; the
emitted vector field and tangency identities are in
\leanmodule{HGAS/Compactification/MixedCentralPolynomialField.lean}.

A delicate point is that chart formulas do not cover the equators
\(w_b=0\).  The module
\leanmodule{HGAS/Compactification/MixedCentralPolynomialGlobalTangency.lean}
proves tangency at \emph{arbitrary} product-sphere points by a direct
polynomial decomposition.  The endpoints
\leanname{eval₂Hom_fieldCode_chartPoint},
\leanname{fieldCode_target_tangent_chartPoint}, and
\leanname{fieldCode_sourceTangentAt_chartPoint} therefore include corners
and equators, not just the finite chart.

\subsection{The compact accumulator}

Given a base drift \(g:B\to\R^m\), the accumulator sphere obeys the
stereographic field above.  Its finite chart records a primitive
\(Z(t)=Z(0)+\int_0^t g(b(s))\,\dd s\).  If one coordinate has nonzero mean
under a base invariant probability, recurrence of a finite accumulator is
impossible: the primitive grows linearly along recurrent generic points.

The general measure obstruction is formalized in
\leanmodule{HGAS/Accumulator/CompactDriftObstruction.lean}.  In particular,
\leanname{finiteChart_measure_eq_zero} proves that invariant mass on the
finite chart vanishes under the nonzero-mean hypothesis.

The reward has three regimes:
\begin{itemize}
\item it is uniformly negative on every old compactification boundary;
\item it is uniformly negative at the accumulator point at infinity;
\item on a designated finite graph it has a fixed positive margin.
\end{itemize}
The abstract sign conclusion is
\leanname{maximalInvariantAverage_neg_iff_of_compactAccumulatorGap} in
\leanmodule{HGAS/Accumulator/CompactSignDichotomy.lean}.

\subsection{Why the sign hypothesis is only ergodic}

An earlier, unnecessarily strong formulation asked for a nonzero drift
coordinate after conditioning \emph{every} invariant probability that gives
positive mass to the finite chart.  That formulation creates a false
burden: different ergodic components can carry opposite drifts and cancel
in a mixture.  The compact-sign proof never needs to rule out such
cancellation.

The corrected theorem uses the ergodic maximizer from
Section~\ref{sec:ergodic-maximizer}.  Let \(\mu_\ast\) maximize the reward
and be ergodic for the full source flow.  The finite chart is measurable and
invariant, so
\[
\mu_\ast(\text{finite chart})\in\{0,1\}.
\]
There are only two cases.
\begin{enumerate}
\item If the mass is zero, the pointwise boundary estimate gives
\(\int\mathcal R\,\dd\mu_\ast\leq-\varepsilon\).
\item If the mass is one, condition on the finite chart and project to the
base.  The compact-accumulator obstruction says that every drift coordinate
has mean zero.  A source-specific theorem saying that a positive finite
conditional of an \emph{ergodic} invariant probability has some nonzero
drift coordinate gives a contradiction.
\end{enumerate}
Thus the maximizing measure must lie in the first case, and the maximal
reward is at most \(-\varepsilon\).

This proof is
\leanname{maximalInvariantAverage_le_neg_epsilon_of_compactAccumulator}.
Its strict two-sided version is
\leanname{maximalInvariantAverage_neg_iff_of_compactAccumulatorGap}.  In
both statements the nonzero-mean hypothesis is quantified only over
\[
\mu:\leanname{InvariantProbabilityMeasure}(\phi)
\quad\text{with}\quad
\leanname{IsErgodicInvariantProbability}(\phi,\mu).
\]
No ergodic decomposition and no no-cancellation theorem for arbitrary
invariant mixtures is assumed.

\section{The concrete compact-accumulator pipeline}
\label{sec:pipeline}

\subsection{Finite code and degree bounds}

The compiler in
\leanmodule{HGAS/Compiler/CompactAccumulatorPipelineCode.lean} combines:
\[
\text{mixed compact base field}
\;+\;\text{slowed drift coordinates}
\;+\;\text{compact accumulator field}.
\]
Its reward is
\[
\mathcal R=W_{\mathrm{acc}}+2W_0-1-\varepsilon,
\]
where \(W_{\mathrm{acc}}\) is the accumulator's last coordinate and
\(W_0=\prod_b w_b^2\) is the product of squared central weights.  On a
finite graph \(W_{\mathrm{acc}}\) is close to \(1\) and \(W_0\) is bounded
below; on a source face \(W_0=0\); at accumulator infinity
\(W_{\mathrm{acc}}=-1\).

Exact chart semantics are the theorems
\leanname{eval₂Hom_fieldCode_base_scaledChart},
\leanname{eval₂Hom_fieldCode_u_scaledChart},
\leanname{eval₂Hom_fieldCode_w_scaledChart}, and
\leanname{eval₂Hom_rewardCode}.  Primitive recursiveness is proved in
\leanmodule{HGAS/Compiler/CompactAccumulatorPipelineComputability.lean}.
The degree module proves the explicit bounds
\[
\deg(\text{field})\leq E+kM+3,\qquad
\deg(\text{reward})\leq\max(1,2k).
\]

\subsection{A concrete fifty-dimensional static linked-source instance}
\label{sec:linked-compact-source}

The generic pipeline has now been instantiated with the linked compiler of
Section~\ref{sec:linked-jacobi-compiler}.  The bookkeeping is completely
explicit:
\[
\begin{array}{c|c}
\text{quantity}&\text{value}\\ \hline
\text{real source coordinates \(n\)}&28\\
\text{real residual drifts \(m\)}&20\\
\text{affine blocks \(k\)}&1\\
\text{already compact source blocks \(\ell\)}&0\\
\text{degree bound and clearing exponent \(E=M\)}&8.
\end{array}
\]
All twenty-eight source coordinates lie in one affine block.  Mixed central
compactification therefore adds one central coordinate, giving a
twenty-nine-coordinate compact base.  Stereographically compactifying the
twenty residual accumulators adds \(20+1=21\) coordinates.  The emitted
field and reward consequently live on
\[
(28+1)+(20+1)=50
\tag{10.1}
\]
real coordinates.

These are literal type-level dimensions in
\leanmodule{HGAS/Modular/LinkedJacobiCompactSourceCode.lean}, not a
paper-only count.  The definitions
\leanname{LinkedJacobiCompactSourceCode.sourceFieldCode} and
\leanname{LinkedJacobiCompactSourceCode.residualCode} insert the exact
twenty-eight-dimensional linked field and its twenty scalar residuals.
The runtime input consists of the two Gaussian-rational clock coefficients,
a rational accumulator scale, and a rational strict-sign margin.
\leanname{LinkedJacobiCompactSourceCode.compile} returns the complete
fifty-dimensional field code together with its reward code.

The two degree ledgers imply the exact affine-block clearing hypotheses via
\leanname{LinkedJacobiCompactSourceCode.sourceFieldCode_termsAffineBlockDegreeLE}
and
\leanname{LinkedJacobiCompactSourceCode.residualCode_termsAffineBlockDegreeLE}.
Specializing the generic bounds gives raw degree at most
\[
8+1\cdot8+3=19
\]
for the emitted field and degree at most \(2\) for the reward (hence the
single combined compiler bound is \(19\)).  The uncompactified source
semantics hold at every point of \(\R^{28}\) by
\leanname{LinkedJacobiCompactSourceCode.complexify_realize_sourceFieldCode}.
Formation of the whole compact field--reward pair is primitive recursive by
\leanname{LinkedJacobiCompactSourceCode.compile_primrec}.

This all-affine choice also resolves one otherwise source-specific premise.
The definition of \leanname{BaseCompactTangency} quantifies only over source
blocks that were compact before clearing.  Here that index type is
\(\operatorname{Fin}(0)\), so
\leanname{LinkedJacobiCompactSourceCode.baseCompactTangency} proves the
premise by empty elimination.  Thus base compact tangency is fully
discharged for this concrete linked block.

The scope is nevertheless important.  This is a proved compact
\emph{compiler instance} for the static unary--Jacobi component.  It does
not adjoin arithmetic detectors or dynamically coupled controller clocks;
it does not construct the invariant component and finite primitive required
by the positive graph; and it does not establish the ergodic nonzero-drift
alternative.  Accordingly it does not inhabit
\leanname{CanonicalCompactAccumulatorGap} or
\leanname{FixedUniversalCompactSignSourceFamily} by itself.

\subsection{A provisional all-affine fifty-eight-dimensional instance}
\label{sec:dynamic-compact-source}

The dynamic compiler of Section~\ref{sec:dynamic-linked-compiler} has also
been inserted literally into the same pipeline.  The module
\leanmodule{HGAS/Modular/TwoControllerLinkedJacobiCompactSourceCode.lean}
fixes
\[
\begin{array}{c|c}
\text{quantity}&\text{value}\\ \hline
\text{real source coordinates \(n\)}&36\\
\text{transported real residuals \(m\)}&20\\
\text{affine blocks \(k\)}&1\\
\text{already compact source blocks \(\ell\)}&0\\
\text{source degree bound}&9\\
\text{residual degree bound}&4\\
\text{even clearing exponent \(M\)}&10.
\end{array}
\tag{10.2}
\]
The twenty residual codes are not reconstructed.  Each is transported from
the rightmost twenty-eight linked variables into the thirty-six-variable
space by the generic right-block embedding, so it ignores exactly the
eight controller--phase coordinates.  The theorems
\leanname{complexifyPoint_eval_residualCode_arbitrary} and
\leanname{forall_eval_residualCode_eq_zero_iff_arbitrary} prove,
respectively, exact evaluation as the ten named complex residuals and the
equivalence between their vanishing and simultaneous vanishing of the
twenty real codes at every ambient point.

All thirty-six source coordinates are assigned to one affine block.  The
source field has raw degree at most \(9\), the residuals have degree at most
\(4\), and both are promoted to the degree-\(10\) affine-block clearing
premise.  The kernel proves that \(10\) is even and \(9<10\).  Since there
are no source blocks declared already compact,
\leanname{TwoControllerLinkedJacobiCompactSourceCode.baseCompactTangency}
again closes \leanname{BaseCompactTangency} by eliminating
\(\operatorname{Fin}(0)\).

The one central source weight and the common accumulator compactification
give the exact type-level output dimension
\[
 (36+1)+(20+1)=58,
\tag{10.3}
\]
proved by
\leanname{TwoControllerLinkedJacobiCompactSourceCode.finalDimension_eq_fifty_eight}.
The generic degree ledger, used with the common promoted source/residual
bound \(E=9\), certifies field degree at most
\[
 9+1\cdot10+3=22
\]
and reward degree at most \(2\).  The complete map from
\((p_\Omega,p_N,\kappa,\varepsilon)\) to the fifty-eight-variable field and
reward is primitive recursive by
\leanname{TwoControllerLinkedJacobiCompactSourceCode.compile_primrec}.

This is a stronger concrete compiler instance than (10.1): the two
controllers are already present and dynamically drive the linked field.
It is also a mathematically valid compactification.  Its drawback is
geometric rather than logical: by placing all thirty-six coordinates in
one affine block, it forgets that the first eight coordinates naturally
form four invariant circles.  Thus it is useful as a simple executable
fallback, but it is no longer the preferred layout for transporting the
clock dynamics.

It is not yet the universal source.  In particular, the arithmetic
detector, the right-hand-side oscillator, their physical factor maps, the
selector component, and the two invariant-measure sign alternatives have
not been composed into this fifty-eight-coordinate object.

\subsection{The geometry-faithful mixed fifty-eight-dimensional instance}
\label{sec:mixed-dynamic-compact-source}

The preferred construction is
\leanmodule{HGAS/Modular/TwoControllerLinkedJacobiMixedCompactSourceCode.lean}.
It uses the canonical split
\[
 \underbrace{\R^8}_{\substack{\text{two controller--phase blocks}\\
                              \text{four compact circles}}}
 \times
 \underbrace{\R^{28}}_{\text{one affine linked-Jacobi block}} .
\]
Within each four-coordinate clock block, \((r,s)\) is one compact circle
and \((c,d)\) is the second.  The two clock blocks therefore give four
already compact source blocks.  Only the rightmost twenty-eight linked
coordinates are centrally compactified.  The reusable layout and its
degree-transfer lemmas are in
\leanmodule{HGAS/Compactification/RightAffineLeftCompactLayout.lean}.

This distinction improves the clearing exponent.  The complete source has
controller/right-linked bidegree at most \((3,8)\), while the residuals have
bidegree at most \((0,4)\).  Denominators arise only from the right affine
variables, so ordinary total degree \(9\) is not the relevant clearing
quantity.  The even exponent
\[
 M=8
\]
already dominates every affine-block degree.  The exact endpoints are
\leanname{TwoControllerLinkedJacobiMixedCompactSourceCode.sourceFieldCode_termsAffineBlockDegreeLE}
and
\leanname{TwoControllerLinkedJacobiMixedCompactSourceCode.residualCode_termsAffineBlockDegreeLE}.

The nonvacuous issue is tangency to the four inherited circles at the
affine equator.  First,
\leanname{TwoControllerLinkedJacobiMixedCompactSourceCode.sourceFieldCode_compactBlockDot}
proves the original
thirty-six-dimensional field tangent to each circle at every ambient
point.  Second, every output coordinate belonging to a compact circle has
right-affine degree \(0\), as recorded by
\leanname{TwoControllerLinkedJacobiMixedCompactSourceCode.sourceFieldCode_termsAffineDegreeZero_of_compact}.
Consequently
denominator clearing multiplies the compact-circle velocity by the common
slowdown without introducing an equator remainder.
\leanname{TwoControllerLinkedJacobiMixedCompactSourceCode.clearedSourceVelocity_compactBlockDot}
proves the resulting dot
product is zero for arbitrary chart coordinates, including \(w=0\), and
\leanname{TwoControllerLinkedJacobiMixedCompactSourceCode.baseCompactTangency}
supplies the pipeline premise globally.  No division by the central
coordinate and no interior-chart assumption occurs.

There is still one affine central coordinate, so the numerical output
dimension is unchanged:
\[
 (36+1)+(20+1)=58.
\]
Geometrically, however, the compact base now has five sphere blocks: four
inherited circles and one sphere compactifying the linked block.  The
accumulator contributes a sixth sphere.  These counts are kernel-checked
by \leanname{TwoControllerLinkedJacobiMixedCompactSourceCode.sphereBlockCounts}.
Using the common total-degree bound \(E=9\), the generic ledger gives field
degree at most
\[
 9+1\cdot8+3=20
\]
and reward degree at most \(2\).  Exact residual semantics, the
fifty-eight-coordinate type identity, and primitive-recursive compilation
are respectively exposed by
\leanname{TwoControllerLinkedJacobiMixedCompactSourceCode.forall_eval_residualCode_eq_zero_iff_arbitrary},
\leanname{TwoControllerLinkedJacobiMixedCompactSourceCode.finalDimension_eq_fifty_eight},
and
\leanname{TwoControllerLinkedJacobiMixedCompactSourceCode.compile_primrec}.

This mixed compiler faithfully preserves the clock product geometry needed
by the analytic Haar theorem.  It still does not construct the complete
semiconjugacy from the compactified polynomial source to the analytic clock
flow, nor does it add the detector, oscillator, or source-specific sign
certificates.

\subsection{Global tangency and boundary reward}

The complete product-sphere tangency theorem is
\leanname{CompactAccumulatorPipelineCode.fieldCode_sourceTangentAt} in
\leanmodule{HGAS/Compiler/CompactAccumulatorPipelineTangency.lean}.
It is a global pointwise result, including all equators and corners.  Its
sole source-specific premise is tangency of those source blocks that were
already compact before mixed-central clearing.  The gap module names this
visible premise \leanname{BaseCompactTangency}; it is not silently produced
by the generic compiler.

The reward bounds are in
\leanmodule{HGAS/Compiler/CompactAccumulatorPipelineReward.lean}:
\leanname{eval₂Hom_rewardCode_le_neg_epsilon_on_sourceFace},
\leanname{eval₂Hom_rewardCode_le_neg_epsilon_at_infinity}, and
\leanname{finiteGraphGap_le_eval₂Hom_rewardCode}.

\subsection{Interior invariance and the base factor}

The file
\leanmodule{HGAS/Compiler/CompactAccumulatorPipelineInterior.lean}
defines the finite interior as the complement of all source equators and
accumulator infinity.  Every boundary component is invariant because its
defining coordinate solves a scalar linear equation \(y'=a(t)y\); uniqueness
forces a zero initial value to remain zero.  Theorems
\leanname{isInvariant_sourceBoundarySet} and
\leanname{isInvariant_sourceFiniteInteriorSet} formalize the union and
complement statements.  The same module proves the pointwise bound
\leanname{sourceRewardBCF_le_neg_epsilon_of_not_mem_finiteInterior}.

On the finite interior, the full flow factors through the compact base.
The module
\leanmodule{HGAS/Compiler/CompactAccumulatorPipelineFactor.lean} constructs:
\begin{itemize}
\item the continuous base projection \leanname{baseProjection};
\item the canonical base flow \leanname{baseFlow};
\item the semiconjugacy \leanname{baseProjection_flow_eq} and its restricted
form \leanname{canonicalRestrictedBaseProjection_semiconj};
\item the finite accumulator projection
\leanname{finiteAccumulatorProjection};
\item the actual bounded continuous base drift \leanname{baseDriftBCF};
\item the \(u\)- and \(w\)-coordinate ODEs
\leanname{canonicalFiniteAccumulator_u_hasDerivAt} and
\leanname{canonicalFiniteAccumulator_w_hasDerivAt}.
\end{itemize}
The semiconjugacy is proved by ODE uniqueness, not by comparing only formal
expressions on a dense chart.

\subsection{Finite graphs on one invariant base component}
\label{sec:component-graph}

A positive finite graph should not be required over the whole compact base.
The canonical no-instance primitive is expected only on one selected
modular/Jacobi component, while unrelated components can meet central
equators.  A global positive lower bound for the central-weight product
would therefore be both unnatural and, in the presence of such equators,
impossible.

The module
\leanmodule{HGAS/Compiler/CompactAccumulatorPipelineComponentGraph.lean}
implements the correct local interface.  A
\leanname{BaseInvariantComponent} consists of a measurable invariant subset
\(C\) of the compact base.  Its subtype \leanname{BaseInvariantComponent.Space}
carries the restricted flow \leanname{BaseInvariantComponent.flow}.  One
then supplies a continuous primitive
\[
z:C\longrightarrow\R^m
\]
whose derivative along the restricted flow is the compiled base drift:
\[
\frac{\dd}{\dd t}z_i(\phi_t b)=g_i(\phi_t b).
\]
No value or extension of \(z\) outside \(C\) is chosen.

Inverse stereographic projection gives the graph
\[
G(b)=\bigl(b,u(z(b)),w(z(b))\bigr)
\]
inside the complete product-sphere source.  The theorem
\leanname{componentFiniteGraphPoint_isIntegralCurve} checks every emitted
polynomial coordinate: the base coordinates use the base flow, and the
accumulator coordinates use the exact derivatives of \(u\) and \(w\).
The theorem \leanname{componentFiniteGraphPoint_flow_eq} then invokes ODE
uniqueness on the closed unit ball to identify this explicit graph
trajectory with the canonical complete source flow.  Equivariance is
packaged as \leanname{componentFiniteGraphPoint_semiconj}.

Consequently any invariant probability \(\eta\) for the restricted flow on
\(C\) pushes through \(G\) to an invariant probability of the \emph{complete}
source; this is \leanname{componentFiniteGraphInvariantProbability}.  The
quantitative hypotheses may also be restricted to this one probability.
If, \(\eta\)-almost everywhere,
\[
\|z(b)\|^2\leq \frac{w_\ast}{4},
\qquad
w_\ast\leq W_0(b),
\]
then
\[
\frac{3w_\ast}{2}-\varepsilon
\leq
\int\mathcal R\,\dd(G_\ast\eta).
\tag{10.4}
\]
The pointwise algebra behind (10.4) uses the inverse-stereographic lower
bound for \(w(z)\) and the reward formula
\(\mathcal R=W_{\rm acc}+2W_0-1-\varepsilon\).  The integrated result is
\leanname{componentFiniteGraphInvariantProbability_reward_lower_ae}; the
existence-shaped endpoint is
\leanname{exists_componentFiniteGraphInvariantProbability_reward_lower_ae}.
Notice that neither inequality is asserted on the rest of the ambient base.

For source assembly, \leanname{ComponentFiniteGraphCertificate} packages
the component, primitive, derivative equation, component invariant
probability, \(w_\ast\), and the two almost-everywhere bounds.  Its endpoint
\leanname{ComponentFiniteGraphCertificate.exists_sourceProbability_reward_lower}
returns exactly the complete-source invariant probability and reward lower
bound needed downstream.  The bridge
\leanname{CanonicalCompactAccumulatorGap.finite_graph_of_componentFiniteGraphCertificate}
converts such a certificate into the positive-graph field of the canonical
gap structure after checking the desired margin \(\delta\).

Every graph, ODE, equivariance, pushforward, and reward theorem in this
subsection is kernel checked.  What remains open is source specific:
construct the canonical component \(C\), its primitive \(z\), its invariant
probability, and the two quantitative bounds from the modular/Jacobi
selector.  The definition of
\leanname{ComponentFiniteGraphCertificate} does not assert that these data
exist.

\subsection{The exact remaining semantic gap}
\label{sec:open-source}

After the separately visible source premise
\leanname{BaseCompactTangency} has been supplied, the remaining hypotheses
are collected in
\leanmodule{HGAS/Compiler/CompactAccumulatorPipelineGapCertificate.lean}.
The structure \leanname{CanonicalCompactAccumulatorGap} has exactly two
source-specific dynamical obligations:

\begin{enumerate}
\item \textbf{Nonzero mean in a solvable instance.}  Every
\emph{ergodic} invariant probability giving positive mass to the finite
interior has, after conditioning and projection to the base, some drift
coordinate of nonzero mean.
\item \textbf{Positive finite graph in an unsolvable instance.}  There is
an invariant probability on the designated finite graph whose reward has a
fixed positive lower bound.
\end{enumerate}

The word ``ergodic'' in the first field is essential and is present in the
actual Lean structure.  Section~\ref{sec:ergodic-maximizer} explains why
this weaker, correct quantifier is sufficient.  Requiring the conclusion
for all invariant mixtures would add an unproved and unnecessary
no-cancellation assertion.  For the second field,
\leanname{finite_graph_of_componentFiniteGraphCertificate} supplies the
measure-level conclusion from the component data of
Section~\ref{sec:component-graph}.

Given these two fields,
\leanname{CanonicalCompactAccumulatorGap.toCompactSignSourceCertificate}
automatically supplies tangency, finite-interior invariance, the factor map,
the accumulator equations, and boundary negativity.  The conditional
endpoint is
\leanname{CanonicalCompactAccumulatorGap.finalField_gas_iff_solvable}.

\statusopen\quad The universal arithmetic construction does not yet supply
an inhabitant of \leanname{CanonicalCompactAccumulatorGap}.  This is where
the all-state modular/Jacobi selector, the required Haar exclusion,
detector realization, physical/additive normalization of the
right-hand-side oscillator, the canonical no-instance component
certificate, and the ergodic nonzero-mean theorem must be assembled.  The
controller part of this measure argument is no longer open:
\leanname{exists_lockedController_sign_and_phaseHaar} supplies the common
locked-sign tuple, the almost-everywhere locked-fibre semiconjugacy, and
the Haar phase marginal for the literal coupled clock flow.  What remains
is to construct the exact factor from the complete compactified
modular/Jacobi/detector source to that flow and transport the conclusion
through the chosen interior component.

The visible base-tangency premise has already been discharged for the
static all-affine twenty-eight-dimensional linked block by
\leanname{LinkedJacobiCompactSourceCode.baseCompactTangency}, and the
provisional all-affine thirty-six-dimensional two-controller block by
\leanname{TwoControllerLinkedJacobiCompactSourceCode.baseCompactTangency}.
More importantly, the geometry-faithful layout with four inherited circles
and one affine linked block has the nonvacuous global theorem
\leanname{TwoControllerLinkedJacobiMixedCompactSourceCode.baseCompactTangency},
including the affine equator.  None of these facts automatically supplies
tangency after detector and oscillator coordinates are added or if the
final source chooses a different enlarged layout.  The general
component-graph and ergodic-maximizer mechanisms are no longer missing.
This source-specific assembly is the first missing link between the fixed
DQL theorem of Section~\ref{sec:arithmetic} and the homogeneous GAS field.

\section{The homogeneous product-sphere wrapper}
\label{sec:wrapper}

The wrapper is complete independently of the open source assembly.
Its input is:
\begin{itemize}
\item a rational polynomial flow tangent to a product of sphere blocks;
\item a rational polynomial reward \(a\);
\item an odd chart degree \(D\);
\item a source flow solving the literal polynomial ODE.
\end{itemize}

\subsection{Two sheets and polynomial charting}

An extra chart coordinate \(z\) is added.  On the normalized unit sphere,
the full stratum has two sheets \(z=\pm1/\sqrt L\), where \(L\) is the
number of blocks.  A source monomial of degree at most \(D\) is homogenized
with powers of \(z\).  Because \(D\) is odd,
\[
L^{(D-1)/2}z^{D-1}=1
\]
on both sheets.  Hence the positive and negative sheets carry the same
source time orientation and the same reward, while all emitted
coefficients remain rational.

The sparse code is built in
\leanmodule{HGAS/Wrapper/UnitSpherePolynomialWrapper.lean}.  The final code
\leanname{finalFieldCode} is proved homogeneous by
\leanname{finalFieldCode_isHomogeneous_of_degreeBounds}; its target degree
is proved odd by \leanname{targetDegree_odd}.

\subsection{Radius cleanup and invariant faces}

Let \(s_i\) be the scaled squared radius of block \(i\).  The cleanup field
is arranged so that a nonnegative dissipation vanishes precisely when each
radius is either \(0\) or \(1\).  Integrating the derivative identity
against an invariant probability forces the dissipation to vanish almost
everywhere.  Thus almost every point belongs either to:
\[
\text{the full stratum }(s_i=1\ \forall i),
\quad\text{or a proper face }(s_i=0\text{ for some }i).
\]
This classification is proved in
\leanmodule{HGAS/Wrapper/RadiusInvariantFaces.lean}.

The wrapper reward equals the source reward on the full stratum and equals
\(-1\) on every proper face.  Splitting an arbitrary invariant probability
between these two invariant pieces gives the exact maximal-average formula
\[
\beta_{\mathrm{wrapper}}=\max\{\beta_{\mathrm{source}},-1\}.
\]
The measure decomposition and attainment proof are in
\leanmodule{HGAS/Wrapper/InvariantAverage.lean} and the concrete sheet
transport is in
\leanmodule{HGAS/Wrapper/UnitSphereWrapperInvariantAverage.lean}.

\subsection{The GAS equivalence}

The positive and negative sheet flows are shown to solve the same source
polynomial ODE.  ODE uniqueness proves that they coincide after pullback;
see \leanname{positiveSourceFlow_isIntegralCurve_sourceField} and
\leanname{positiveSourceFlow_eq_negativeSourceFlow} in
\leanmodule{HGAS/Wrapper/UnitSphereWrapperSourceFlow.lean}.

The final concrete theorem is
\[
\boxed{
\text{the wrapped homogeneous polynomial field is GAS}
\Longleftrightarrow
\text{the source maximal invariant reward is negative}.}
\]
Its Lean name is
\leanname{realize_finalFieldCode_gas_iff_externalSourceMaximum_neg} in
\leanmodule{HGAS/Wrapper/UnitSphereWrapperConcreteReduction.lean}.

The compiler-facing structure
\leanname{CompactSignSourceCertificate} and its endpoint
\leanname{CompactSignSourceCertificate.finalField_gas_iff_solvable} occur
in \leanmodule{HGAS/Compiler/CompactSignSourceWrapper.lean}.  Once the open
gap certificate of Section~\ref{sec:open-source} is constructed, no
wrapper-side dynamics or measure theory remains.

\section{Effective cubicization}
\label{sec:cubic}

Suppose \(F:\R^n\to\R^n\) is homogeneous of odd degree \(2k+1\).  Let
\[
\nu_k(x)=(x^\alpha)_{|\alpha|=k}
\]
be the degree-\(k\) Veronese map.  Its coordinates satisfy finitely many
quadratic toric relations.  Differentiating \(\nu_k(x(t))\) along
\(\dot x=F(x)\) produces expressions of degree \(3k\), which can be
factored into cubic polynomials in the variables \(y_\alpha=x^\alpha\).
This gives a cubic vector field \(G\) satisfying
\[
D\nu_k(x)\,F(x)=G(\nu_k(x)).
\]
This construction is logically separate from earlier complexity results
for cubic homogeneous stability, such as Ahmadi~\cite{Ahmadi2012}: here the
task is an exact GAS-preserving compiler from every fixed odd degree to
degree three.

\subsection{Geometry and transverse stabilization}

The real Veronese cone is the common zero set of explicit quadratic
relations.  Merely extending the source dynamics to an ambient cubic
field is not enough: trajectories away from the cone could spoil GAS.
The construction therefore adds a cubic negative gradient of the squared
quadratic residuals.  Near the cone this supplies normal contraction; away
from it, compactness gives a uniform penalty.

The geometric development occupies \leanmodule{HGAS/Veronese}.  The
central semantic equivalence is
\leanname{sourceGAS_iff_stabilizedCubicGAS} in
\leanmodule{HGAS/Veronese/CubicLiftGASConjugacy.lean}.  Its proof uses:
\begin{itemize}
\item exact conjugacy on the real cone;
\item a Morse--Bott-type coercive estimate transverse to the cone;
\item an invariant-measure sandwich showing that invariant mass outside the
controlled cone tube vanishes;
\item the maximal-angular-average GAS criterion of Section~\ref{sec:problem}.
\end{itemize}

The algebraic background is classical Veronese geometry; see
Hartshorne~\cite{Hartshorne1977} and Sturmfels~\cite{Sturmfels1996}.  The
formal development proves the particular finite toric and differential
identities needed here.

\subsection{A computable stabilizing gain}

A semantic compactness proof could choose a real gain nonconstructively,
which would be useless for a many-one reduction.  The effective module
\leanmodule{HGAS/Computability/CubicLiftEffectiveGain.lean} instead bounds
all source-dependent terms linearly by the coefficient \(\ell^1\)-mass.
Geometry-dependent constants are frozen for fixed \(n,k\), and the input
code supplies an explicit rational majorant.

The gain code \leanname{effectiveCubicGainCode} is primitive recursive and
positive.  The total compiler \leanname{effectiveCubicReductionCode}
includes a malformed-input guard, emits a syntactically homogeneous cubic
field for every raw input, and satisfies
\[
\leanname{HGASQ n (2*k+1) F}
\Longleftrightarrow
\leanname{HGASQ N3 3 (effectiveCubicReductionCode F)}.
\]
The exact theorems are
\leanname{effectiveCubicReductionCode_hgasQ_iff} and
\leanname{hgasQ_manyOneReducible_cubic}.  Sparse polynomial realization is
verified in
\leanmodule{HGAS/Computability/CubicLiftCodeCompiler.lean}.

The effectivity statement is fixed-parameter, exactly as required for a
fixed-language many-one reduction.  The geometric natural-number multiplier
in \leanname{fixedCubicGainParameters} is selected nonuniformly once
\((n,k)\) is fixed.  It is then a hard-coded constant, and the resulting map
of input field codes is proved computable.  The present theorem does not
compute a numerical multiplier uniformly from varying \(n,k\).

\section{Dimension and degree propagation}
\label{sec:propagation}

\subsection{Adding dimensions}

To add \(p\) coordinates, take the direct product of a cubic field \(G(y)\)
with independent cubic sinks
\[
\dot z_j=-z_j^3.
\]
The sink is forward complete, stable, and globally attractive.  The direct
product is GAS exactly when \(G\) is GAS.

\subsection{Raising odd degree}

For \(j\geq0\), multiply the entire product field by
\[
\bigl(\|y\|^2+\|z\|^2\bigr)^j.
\]
Away from the origin this is a positive state-dependent time change, so it
does not change orbits or their limiting behavior.  At the origin the
homogeneous extension is still a polynomial field.  Therefore a cubic
field is GAS exactly when the resulting field of degree \(3+2j\) is GAS.
The semantic statement is \leanname{degreePaddedField_gas_iff} in
\leanmodule{HGAS/Propagation/DegreePaddingGAS.lean}.

\subsection{The finite compiler}

The code module
\leanmodule{HGAS/Propagation/DegreePaddingCode.lean} proves exact coordinate
insertion, decoding, homogeneity, realization, primitive recursiveness, and
the GAS equivalence.  Combining it with cubicization gives
\leanname{oddToCofinalCode}.  For all raw inputs,
\[
\HGAS_{\Q}(n,2k+1;F)
\Longleftrightarrow
\HGAS_{\Q}\bigl(N, d;
  \leanname{oddToCofinalCode}(F)\bigr)
\]
whenever the target dimension is large enough and \(d\geq3\) is odd.

The principal many-one theorem is
\leanname{hgasQ_manyOneReducible_cofinalFixedPair} in
\leanmodule{HGAS/Propagation/CofinalFixedPairReduction.lean}.  The final
recursion-theoretic composition
\leanname{sigmaOneComplete_cofinalFixedPair_of_universalHalts_reduction}
is in \leanmodule{HGAS/Propagation/CofinalSigmaOne.lean}.

Thus cubicization and cofinal propagation are not open mathematical
steps.  What is missing is the upstream fixed-pair reduction from universal
halting to one homogeneous GAS language.

\section{What a completed headline theorem will look like}
\label{sec:headline}

The structure
\leanname{FixedUniversalCompactSignSourceFamily} in
\leanmodule{HGAS/Compiler/UnconditionalFixedCompactSignSourceFamily.lean}
is the precise final source-assembly target.  It asks for:
\begin{itemize}
\item one fixed product-sphere layout;
\item one fixed odd wrapper degree;
\item for every universal-halting input, a
\leanname{CompactSignSourceCertificate};
\item a computable finite map producing the corresponding field and reward
codes;
\item equality between those finite codes and the semantic certificate data.
\end{itemize}

Importantly, the file makes no existence claim.  If a family is supplied,
the already proved theorem \leanname{fixedPair_sigmaOneComplete} gives
\(\Sigma^0_1\)-completeness at one fixed dimension and odd degree.  The
already proved theorem \leanname{cofinal_sigmaOneComplete} then gives
\(\Sigma^0_1\)-completeness for every sufficiently large fixed dimension
and every fixed odd degree at least three.

The remaining work can therefore be stated cleanly.

\begin{theorem}[Open integration theorem]
Construct a fixed product-sphere layout, a fixed odd wrapper degree, and an
inhabitant of \leanname{FixedUniversalCompactSignSourceFamily} by compiling
the unconditional fixed DQL family through the modular/Jacobi selector and
the canonical compact-accumulator gap.  Starting from the proved
thirty-six-coordinate dynamic linked source and its geometry-faithful
fifty-eight-coordinate mixed compact pipeline, integrate the physical
detector and right-hand-side oscillator, construct the exact factor from
the enlarged source to the already classified coupled controller--phase
flow, and identify its recurrent modular factor with the Jacobi selector.
Then construct the remaining
\leanname{BaseCompactTangency} proof for the final combined layout, the
\leanname{ComponentFiniteGraphCertificate} in the no-instance case, and the
corrected ergodic nonzero-mean alternative in the solvable case.
\end{theorem}

Once this theorem is kernel checked, the final proof is a short composition:
\[
\begin{aligned}
\text{universal halting}
&\leq_m \text{fixed DQL solvability}\\
&\leq_m \HGAS_{\Q}(N_\ast,d_\ast)
\leq_m \HGAS_{\Q}(N,3)
\leq_m \HGAS_{\Q}(N',d),
\end{aligned}
\]
combined with the c.e. upper bound of Section~\ref{sec:upper}.  Positive
denominator clearing will give the corresponding integer-coefficient
statement without changing trajectories.

\section{Module guide}

\begin{center}
\small
\begin{tabular}{>{\raggedright\arraybackslash}p{0.29\textwidth}
                >{\raggedright\arraybackslash}p{0.48\textwidth}
                >{\raggedright\arraybackslash}p{0.13\textwidth}}
\toprule
Mathematical layer & Principal module or endpoint & Status\\
\midrule
GAS definition &
\leanmodule{Dynamics/Stability.lean} & \statusproved\\
Angular maximal-average criterion &
\leanname{globallyAsymptoticallyStable_iff_maximalAngularAverage_neg}
& \statusproved\\
Full-flow ergodic maximizer &
\leanname{exists_ergodicInvariantProbability_isMaxOn_average}
& \statusproved\\
Odd-degree c.e. upper bound &
\leanname{hgasQ_re_of_oddDegree} & \statusproved\\
Primitive-recursive MRDP &
\leanname{binaryPrimitiveRecursiveRelationsAreDiophantine}
& \statusproved\\
Fixed DQL completeness &
\leanname{unconditionalFixedDQLSolvability_sigmaOneComplete}
& \statusproved\\
Controller and radical clocks &
\leanmodule{Clocks/ControllerFlow.lean},
\leanmodule{Clocks/RadicalClock.lean} & \statusproved\\
Radians/turns normalization and physical phase Haar &
\leanname{unitAddCircleHomeomorph_semiconj},
\leanname{physicalSignedRadicalPhase_equivariantMarginal_eq_haar}
& \statusproved\\
Fixed-support and controller-atom selection &
\leanname{exists_eq_diracProba_of_isErgodicInvariantProbability_ae_fixed},
\leanname{exists_eq_diracProba_at_controllerEquilibrium}
& \statusproved\\
Exact coupled controller--phase flow &
\leanname{controllerPhaseFlow_solves_controllerField},
\leanname{controllerPhaseFlow_solves_lockedPhaseField}
& \statusproved\\
Controller polynomial compiler &
\leanname{ControllerPolynomialCode.fieldCode_primrec},
\leanname{ControllerPolynomialCode.fieldCode_termsDegreeLE}
& \statusproved\\
Physical-clock detector average &
\leanname{maximalInvariantAverage_unconditionalUniversalHalting_physicalClocks_iff}
& \statusproved\\
RHS oscillator polynomial compiler &
\leanname{RHSOscillatorPolynomialCode.realize_fieldCode},
\leanname{RHSOscillatorPolynomialCode.fieldCode_primrec}
& \statusproved\\
RHS oscillator graph classification &
\leanname{rhsErgodicInvariantProbability_graphClassification},
\leanname{rationalRhsErgodicInvariantProbability_graphClassification}
& \statusproved\\
Full coupled locked-sign/Haar composition &
\leanname{exists_lockedController_sign_and_phaseHaar},
\leanname{phaseProjection_ae_semiconj} & \statusproved\\
Theta modular identities &
\leanmodule{Modular/ThetaDerivativeIdentities.lean},
\leanmodule{Modular/RamanujanIdentities.lean} & \statusproved\\
Jacobi Wronskians and Delta identity &
\leanname{canonicalJacobiThetaWronskiansAt},
\leanname{jacobiThetaDeltaIdentity} & \statusproved\\
Ten-coordinate Jacobi compiler &
\leanmodule{Modular/JacobiPolynomialCode.lean},
\leanname{JacobiPolynomialCode.compileFieldCode_primrec} & \statusproved\\
Arbitrary complex-family realification &
\leanname{ComplexPolynomialFamilyCode.forall_eval_realify_eq_zero_iff},
\leanname{ComplexPolynomialFamilyCode.realify_primrec} & \statusproved\\
Linked unary--Jacobi compiler &
\leanname{LinkedJacobiPolynomialCode.compileFieldCode_primrec},
\leanname{LinkedJacobiPolynomialCode.forall_eval_realResidualCode_eq_zero_iff}
& \statusproved\\
Dynamic 36-coordinate two-controller linked compiler &
\leanname{TwoControllerLinkedJacobiPolynomialCode.realize_fieldCode_arbitrary},
\leanname{TwoControllerLinkedJacobiPolynomialCode.fieldCode_termsDegreeLE}
& \statusproved\\
Projective Jacobi saturation &
\leanname{range_projectiveJacobiTheta_eq_residualCurve}
& \statusproved\\
Horocycle classification &
\leanname{GammaTwoDaniSmillieTheorem} & \statusopen\\
Conditional recurrent marginal classification &
\leanname{GammaTwoDaniSmillieTheorem.classifyErgodicRecurrentStateMarginal}
& \statusinterface\\
All-state selector/state map &
No completed endpoint & \statusopen\\
Mixed compactification and pipeline factor &
\leanmodule{Compiler/CompactAccumulatorPipelineFactor.lean}
& \statusproved\\
Concrete 50-dimensional linked pipeline &
\leanname{LinkedJacobiCompactSourceCode.compile_primrec},
\leanname{LinkedJacobiCompactSourceCode.baseCompactTangency}
& \statusproved\\
Provisional all-affine 58-dimensional pipeline &
\leanname{TwoControllerLinkedJacobiCompactSourceCode.compile_primrec},
\leanname{TwoControllerLinkedJacobiCompactSourceCode.finalDimension_eq_fifty_eight}
& \statusproved\\
Geometry-faithful mixed 58-dimensional pipeline &
\leanname{TwoControllerLinkedJacobiMixedCompactSourceCode.baseCompactTangency},
\leanname{TwoControllerLinkedJacobiMixedCompactSourceCode.compile_primrec}
& \statusproved\\
Ergodic-only compact sign theorem &
\leanname{maximalInvariantAverage_neg_iff_of_compactAccumulatorGap}
& \statusproved\\
Component-restricted finite graph &
\leanname{ComponentFiniteGraphCertificate.exists_sourceProbability_reward_lower}
& \statusproved\\
Arithmetic compact gap &
\leanname{CanonicalCompactAccumulatorGap} & \statusinterface\\
Product-sphere wrapper &
\leanname{realize_finalFieldCode_gas_iff_externalSourceMaximum_neg}
& \statusproved\\
Fixed universal source family &
\leanname{FixedUniversalCompactSignSourceFamily} & \statusinterface\\
Effective cubicization &
\leanname{hgasQ_manyOneReducible_cubic} & \statusproved\\
Cofinal propagation &
\leanname{hgasQ_manyOneReducible_cofinalFixedPair} & \statusproved\\
\bottomrule
\end{tabular}
\end{center}

\section{Concluding perspective}

The formalization has already changed the proof in useful ways.  It forces
forward completeness to be explicit in the definition of GAS; separates
finite compilers from noncomputable semantics; replaces vague
``generic-orbit'' claims by invariant-probability statements; proves that a
maximal invariant average is realized by a probability ergodic for the full
real flow; removes an unnecessary no-cancellation quantifier from the
compact-sign argument; localizes the no-instance finite graph to its actual
invariant base component; handles compactification equators and wrapper
faces globally; checks every theta normalization; makes both the
ten-coordinate Jacobi field and the linked fourteen-complex-coordinate
field effective; provides an arbitrary-family complex realifier;
instantiates the linked block as a fifty-dimensional static compact
compiler; constructs the dynamically controlled thirty-six-coordinate
source; retains the valid fifty-eight-dimensional all-affine compiler as a
provisional fallback; and constructs the preferred mixed compiler with
four inherited circles, one affine linked block, clearing exponent \(8\),
five base spheres, and global equator tangency.  This keeps the essential
distinction between affine linked degree \(8\) and ordinary total degree
\(9\).  It also makes the radians-to-turns factor \(1/(2\pi)\) explicit,
proves a global coupled controller--phase flow and its ODE, and proves that
every ergodic invariant probability is concentrated on one locked
controller fibre where phase projection is almost everywhere a constant
radical rotation with Haar marginal.  It compiles the factor-four
right-hand-side oscillator, classifies its
ergodic additive-torus measures as Haar on unique graphs, and discharges
the detector's Haar hypothesis whenever the required physical phase factor
is supplied.  Finally, it transports the still-open horocycle
classification, conditionally, through state factors and globally trivial
finite sheets.  It thereby reduces the remaining lower-bound work to named,
mathematically meaningful interfaces.

The present definitive status is therefore:
\begin{quote}
\emph{The c.e. upper bound and all major algebraic, analytic, compiler,
wrapper, ergodic-maximizer, component-graph, cubicization, and propagation
components listed above are kernel checked.  The finite ten-coordinate
Jacobi compiler, linked unary--Jacobi compiler, dynamic two-controller
coupling, both fifty-eight-dimensional compact pipelines, the
geometry-faithful mixed-layout tangency theorem, physical phase
normalization, the locked-fibre controller--phase Haar classification,
oscillator compiler and graph classification, and physical-clock
detector-average implication are complete in their stated scopes.
Conditional on an explicit Dani--Smillie hypothesis and an explicit
modular factor map, the ergodic marginal and recurrence classification is
also complete.

The full \(\Sigma^0_1\)-hardness theorem for homogeneous polynomial GAS is
not yet a theorem of the Lean development.  It still requires: a proof of
the stated Dani--Smillie proposition in the exact quotient; the recurrent
modular/Jacobi state map and selector, including any nontrivial finite-cover
transport; an exact factor from the geometry-faithful compact polynomial
source to the already classified coupled clock flow; exact identification
of the physical polynomial oscillator and detector with their additive
factor models; and construction of the canonical fixed source family,
including tangency for any further enlarged final layout, the
source-specific positive component certificate, and the ergodic
nonzero-drift alternative.  None of those remaining integration
statements, and hence no full lower bound, is claimed here.}
\end{quote}

\begin{thebibliography}{99}
\small

\bibitem{Ahlfors1966}
L.~V. Ahlfors,
\newblock \emph{Complex Analysis: An Introduction to the Theory of Analytic
Functions of One Complex Variable}, 2nd ed.,
\newblock McGraw--Hill, New York, 1966, pp.~271--274.

\bibitem{Ahmadi2012}
A.~A. Ahmadi,
\newblock On the difficulty of deciding asymptotic stability of cubic
homogeneous vector fields,
\newblock in \emph{2012 American Control Conference}, 2012,
pp.~3334--3339,
\newblock \href{https://doi.org/10.1109/ACC.2012.6315687}
{doi:10.1109/ACC.2012.6315687}.

\bibitem{AhmadiElKhadir2020}
A.~A. Ahmadi and B. El Khadir,
\newblock On algebraic proofs of stability for homogeneous vector fields,
\newblock \emph{IEEE Transactions on Automatic Control} 65 (2020), no.~1,
325--332,
\newblock \href{https://doi.org/10.1109/TAC.2019.2914968}
{doi:10.1109/TAC.2019.2914968}.

\bibitem{Brezhnev2013}
Y.~V. Brezhnev,
\newblock Non-canonical extension of theta-functions and modular
integrability of theta-constants,
\newblock \emph{Proceedings of the Royal Society of Edinburgh Section A:
Mathematics} 143 (2013), no.~4, 689--738,
\newblock \href{https://doi.org/10.1017/S0308210512001023}
{doi:10.1017/S0308210512001023}.

\bibitem{BrezhnevLyakhovichSharapov2011}
Y.~V. Brezhnev, S.~L. Lyakhovich, and A.~A. Sharapov,
\newblock Dynamical systems defining Jacobi's theta-constants,
\newblock \emph{Journal of Mathematical Physics} 52 (2011), no.~11,
112704,
\newblock \href{https://doi.org/10.1063/1.3662961}
{doi:10.1063/1.3662961}.

\bibitem{Chandrasekharan1985}
K. Chandrasekharan,
\newblock \emph{Elliptic Functions},
\newblock Grundlehren der mathematischen Wissenschaften 281,
Springer, Berlin, 1985,
\newblock \href{https://doi.org/10.1007/978-3-642-52244-4}
{doi:10.1007/978-3-642-52244-4}.

\bibitem{DLMF}
F.~W.~J. Olver, A.~B. Olde Daalhuis, D.~W. Lozier, B.~I. Schneider,
R.~F. Boisvert, C.~W. Clark, B.~R. Miller, B.~V. Saunders, H.~S. Cohl,
and M.~A. McClain (eds.),
\newblock \emph{NIST Digital Library of Mathematical Functions},
\newblock \url{https://dlmf.nist.gov/}; Sections 20.2, 20.7, and 23.18.

\bibitem{DaniSmillie1984}
S.~G. Dani and J. Smillie,
\newblock Uniform distribution of horocycle orbits for Fuchsian groups,
\newblock \emph{Duke Mathematical Journal} 51 (1984), no.~1, 185--194,
\newblock \href{https://doi.org/10.1215/S0012-7094-84-05110-X}
{doi:10.1215/S0012-7094-84-05110-X}.

\bibitem{DavisPutnamRobinson1961}
M. Davis, H. Putnam, and J. Robinson,
\newblock The decision problem for exponential Diophantine equations,
\newblock \emph{Annals of Mathematics}, Second Series, 74 (1961), no.~3,
425--436,
\newblock \href{https://doi.org/10.2307/1970289}
{doi:10.2307/1970289}.

\bibitem{Hartshorne1977}
R. Hartshorne,
\newblock \emph{Algebraic Geometry},
\newblock Graduate Texts in Mathematics 52, Springer, New York, 1977,
\newblock \href{https://doi.org/10.1007/978-1-4757-3849-0}
{doi:10.1007/978-1-4757-3849-0}.

\bibitem{HubbardMiller2010}
J.~H. Hubbard and R.~L. Miller,
\newblock Equidistribution of horocyclic flows on complete hyperbolic
surfaces of finite area,
\newblock in \emph{In the Tradition of Ahlfors--Bers, V},
Contemporary Mathematics 510, American Mathematical Society, Providence,
2010, pp.~129--150,
\newblock \href{https://doi.org/10.1090/conm/510/10022}
{doi:10.1090/conm/510/10022}.

\bibitem{Lagrange1770}
J.-L. Lagrange,
\newblock Demonstration d'un theoreme d'arithmetique,
\newblock \emph{Nouveaux Memoires de l'Academie Royale des Sciences et
Belles-Lettres de Berlin}, volume for 1770 (printed 1772), 123--133;
reprinted in \emph{Oeuvres de Lagrange}, vol.~3, 1869, pp.~189--201.

\bibitem{Matiyasevich1970}
Y.~V. Matiyasevich,
\newblock Enumerable sets are Diophantine,
\newblock \emph{Soviet Mathematics Doklady} 11 (1970), no.~2, 354--357.

\bibitem{Matiyasevich1971}
Y.~V. Matiyasevich,
\newblock Diophantine representation of enumerable predicates,
\newblock \emph{Mathematics of the USSR-Izvestiya} 5 (1971), no.~1,
1--28,
\newblock \href{https://doi.org/10.1070/IM1971v005n01ABEH001004}
{doi:10.1070/IM1971v005n01ABEH001004}.

\bibitem{Matiyasevich1993}
Y.~V. Matiyasevich,
\newblock \emph{Hilbert's Tenth Problem},
\newblock Foundations of Computing, MIT Press, Cambridge, MA, 1993.

\bibitem{Mumford1983}
D. Mumford,
\newblock \emph{Tata Lectures on Theta I},
\newblock Progress in Mathematics 28, Birkhauser, Boston, 1983,
\newblock \href{https://doi.org/10.1007/978-1-4899-2843-6}
{doi:10.1007/978-1-4899-2843-6}.

\bibitem{Rosier1992}
L. Rosier,
\newblock Homogeneous Lyapunov function for homogeneous continuous vector
field,
\newblock \emph{Systems \& Control Letters} 19 (1992), no.~6, 467--473,
\newblock \href{https://doi.org/10.1016/0167-6911(92)90078-7}
{doi:10.1016/0167-6911(92)90078-7}.

\bibitem{Sturmfels1996}
B. Sturmfels,
\newblock \emph{Grobner Bases and Convex Polytopes},
\newblock University Lecture Series 8, American Mathematical Society,
Providence, RI, 1996,
\newblock \href{https://doi.org/10.1090/ulect/008}
{doi:10.1090/ulect/008}.

\bibitem{Walters1982}
P. Walters,
\newblock \emph{An Introduction to Ergodic Theory},
\newblock Graduate Texts in Mathematics 79, Springer, New York, 1982,
\newblock \href{https://doi.org/10.1007/978-1-4612-5775-2}
{doi:10.1007/978-1-4612-5775-2}.

\end{thebibliography}

\end{document}
